\chapter{Morphismes lisses: g\'en\'eralit\'es,~propri\'et\'es~diff\'erentielles}
\label{II}
Les renvois \`a l'expos\'e~\ref{I} sont indiqu\'es par~I. On
rappelle que les anneaux sont noeth\'eriens, et les
pr\'esch\'emas localement noeth\'eriens.

\section{G\'en\'eralit\'es}
\label{II.1}

Soit $Y$ un pr\'esch\'ema, soient $t_1,\dots,t_n$ des
ind\'etermin\'ees, on pose
\begin{equation}
\label{eq:II.1.1}
Y[t_1,\dots,t_n]=Y\otimes_\ZZ \ZZ[t_1,\dots,t_n]\quoi.
\end{equation}
Donc $Y[t_1,\dots,t_n]$ est un $Y$-sch\'ema, affine au-dessus
de~$Y$, d\'efini par le faisceau quasi-coh\'erent d'alg\`ebres
$\cal{O}_Y[t_1,\dots,t_n]$. La donn\'ee d'une section de ce
pr\'esch\'ema au-dessus de~$Y$ \'equivaut donc \`a la
donn\'ee de~$n$ sections de~$\cal{O}_Y$ (correspondant aux images
des~$t_i$ par l'homomorphisme correspondant). Si $Y^\prime$ est
au-dessus de~$Y$, on a
\begin{equation}
\label{eq:II.1.2}
Y[t_1,\dots,t_n]\times_Y Y'=Y'[t_1,\dots,t_n]\quoi,
\end{equation}
(ce qui implique que la donn\'ee d'un $Y$-morphisme de~$Y^\prime$
dans $Y[t_1,\dots,t_n]$ \'equivaut \`a la donn\'ee de~$n$
sections de~$\cal{O}_{Y^\prime}$), d'autre part on a
\begin{equation}
\label{eq:II.1.3}
\big(Y[t_1,\dots,t_n]\big)[t_{n+1},\dots,t_m]=Y[t_1,\dots,t_m]\quoi,
\end{equation}
en vertu de la formule analogue pour les anneaux de polyn\^omes
sur~$\ZZ$. La formule~\eqref{eq:II.1.2} implique que
$Y[t_1,\dots,t_n]$ varie fonctoriellement avec~$Y$.

$Y[t_1,\dots,t_n]$ est de type fini et plat au-dessus de~$Y$.

\begin{definition}
\label{II.1.1}
Soit $f\colon X \to Y$ un morphisme, faisant de~$X$ un
$Y$-pr\'esch\'ema. On dit que $f$ est
\emph{lisse}\footnote{Ancienne terminologie: $f$ est \emph{simple}
\index{simple: synonyme d\'esuet de lisse|hyperpage}en~$x$, ou $x$ est un point \emph{simple} pour~$f$. Cette
terminologie pr\^etait \`a confusion dans divers contextes
(alg\`ebres simples, groupes simples) et a d\^u \^etre
abandonn\'ee.}
\index{lisse (morphisme, alg\`ebre)|hyperpage}en $x \in X$, ou que $X$ est \emph{lisse sur~$Y$ en~$x$}, s'il existe
un entier $n \geq 0$, un voisinage ouvert $U$ de~$x$, et un
$Y$-morphisme \'etale de~$U$ dans $Y[t_1,\dots,t_n]$. On dit
que~$f$
(\resp $X$)
est \emph{lisse} s'il est lisse en tous les points de~$X$. Une
alg\`ebre~$B$ sur un anneau~$A$ est dite lisse en un id\'eal
premier $\goth{p}$ de~$B$, si $\Spec(B)$ est lisse sur~$\Spec(A)$
au
point~$\goth{p}$;
$B$ est dite lisse sur~$A$ si elle est lisse sur~$A$ en tout id\'eal
premier~$\goth{p}$ de~$B$. Enfin, un homomorphisme local $A \to B$
d'anneaux locaux est dit lisse (ou $B$ est dite lisse sur~$A$)\footnote{Il vaut mieux dire alors, comme dans EGA IV 18.6.1, que $B$
est \og \emph{essentiellement lisse}\fg sur~$A$.}
\index{essentiellement lisse|hyperpage}si $B$ est localis\'ee d'une alg\`ebre de type fini $B_1$ lisse
sur~$A$.
\end{definition}

On note que la notion de lissit\'e de~$X$ sur~$Y$ est locale sur~$X$
et sur~$Y$; si $X$ est lisse sur~$Y$, il est localement de type fini
sur~$Y$.

\setcounter{subsection}{0}
\begin{proposition}
\label{prop:II.1.1}
L'ensemble des points $x$ de~$X$ en lesquels $f$ est lisse est ouvert.
\end{proposition}

C'est trivial sur la d\'efinition.

\begin{corollaire}
\label{II.1.2}
Si $B$ est lisse sur~$A$ en~$\goth{p}$, alors il est lisse sur~$A$
en~$\goth{q}$ pour tout id\'eal premier~$\goth{q}$ de~$B$ contenu
dans~$\goth{p}$.
\end{corollaire}

\ref{prop:II.1.1} implique aussi que les deux derni\`eres
d\'efinitions~\ref{II.1.1} co\"incident dans leur domaine commun
d'existence.

\begin{proposition}
\label{II.1.3}
\textup{(i)}~Un morphisme \'etale, en particulier une immersion
ouverte, un morphisme identique, est lisse. \textup{(ii)}~Une extension
de la base dans un morphisme lisse donne un morphisme
lisse. \textup{(iii)}~Le compos\'e de deux morphismes lisses est
lisse.
\end{proposition}

(i)~est trivial sur la d\'efinition, on a plus pr\'ecis\'ement:

\begin{corollaire}
\label{II.1.4}
\'etale $=$ quasi-fini $+$ lisse.
\end{corollaire}

(ii)~r\'esulte aussit\^ot du fait analogue pour les morphismes
\'etales (I~\ref{I.4.6}) et pour les projections $Y[t_1,\dots,t_n]
\to Y$ (\cf \eqref{eq:II.1.2}). Pour (iii), cela r\'esulte
formellement du fait que c'est vrai s\'epar\'ement pour
\og \'etale\fg (I~\ref{I.4.6}) et des projections du type
$Y[t_1,\dots,t_n]\to Y$
(\cf \eqref{eq:II.1.3}), et des deux faits cit\'es pour~(ii):
Supposons $Y$ lisse sur~$Z$ et $X$ lisse sur~$Y$, prouvons que $X$ est
lisse sur~$Z$; on peut supposer $Y$ \'etale sur~$Z[t_1,\dots,t_n]$
et $X$ \'etale sur~$Y[s_1,\dots,s_m]$, la premi\`ere
hypoth\`ese implique donc que $Y[s_1,\dots,s_m]$ est \'etale sur~$Z[t_1,\dots,t_n][s_1,\dots,s_m] = Z[t_1,\dots,s_m]$, donc $X$ est
\'etale sur~$Z[t_1,\dots,s_m]$, cqfd.

\begin{remarque}
\label{II.1.5}
L'entier $n$ qui figure dans d\'ef.\, \ref{II.1.1} est bien
d\'etermin\'e, car on constate
aussit\^ot que c'est la dimension de l'anneau local de~$x$ dans sa
fibre $f^{-1}\big(f(x)\big)$. On l'appelle \og dimension relative\fg
de~$X$ sur~$Y$. Elle se comporte additivement pour la composition des
morphismes.
\end{remarque}

\section{Quelques crit\`eres de lissit\'e d'un morphisme}
\label{II.2}

\begin{theoreme}
\label{II.2.1}
Soit $f\colon X\to Y$ un morphisme localement de type fini, soit $x\in
X$ et $y=f(x)$. Pour que $f$ soit lisse en $x$, il faut et il suffit
que \textup{(a)} $f$ soit plat en $x$, et \textup{(b)} $f^{-1}(y)$ soit lisse sur~$\kres(y)$
en $x$.
\end{theoreme}

Le compos\'e de deux morphismes plats \'etant plat, et
$Y[t_{1},\dots,t_{n}]\to Y$ \'etant un morphisme plat, on voit que
lisse implique plat; compte tenu de~\ref{II.1.3} (ii) cela prouve la
n\'ecessit\'e. Supposons (a) et (b) v\'erifi\'ees, soient $V$
un voisinage affine de~$y$ d'anneau~$A$, $U$ un voisinage affine de
$x$ au-dessus de~$V$, d'anneau $B$. Prenant $U$ assez petit, on peut
supposer par (b) qu'il existe un $\kres(y)$-morphisme \emph{étale}
$$
g\colon U|f^{-1}(y)\to\Spec k[t_{1},\dots,t_{n}]\quad (k=\kres(y))
$$
d\'efini par $n$ sections $g_{i}$ du faisceau structural de
$U|f^{-1}(y)$. On constate facilement qu'on peut supposer que les
$g_{i}$ (qui a priori sont des \'el\'ements de
$B\otimes_{A}k=BS^{-1}$, o\`u $S=A-\goth{p}$, $\goth{p}$ l'id\'eal
premier de~$A$ correspondant \`a $y$) proviennent de sections du
faisceau structural de
$U$,
donc que $g$ est induit par un morphisme,
encore not\'e~$g$
$$
g\colon U\to Y[t_{1},\dots,t_{n}]
$$
(quitte \`a multiplier les $g_{i}$ par un m\^eme \'el\'ement
non nul de~$k$). Or
$U$
est plat sur~$Y$ par (a), il en est de m\^eme de
$Y[t_{1},\dots,t_{n}]$, d'autre part $g$ induit un morphisme
\'etale entre les fibres au-dessus de~$y$, donc $g$ est \'etale en
$x$ par~(I~\ref{I.5.8}), cqfd.

\begin{corollaire}
\label{II.2.2}
Soient $S$ un pr\'esch\'ema, $f\colon X\to Y$ un $S$-morphisme de
type fini, $Y$ \'etant de type fini et plat sur~$S$, $x\in X$, $s$
la projection de~$x$ sur~$S$. Pour que $f$ soit lisse en $x$, il faut
et il suffit que $X$ soit plat (ou encore: lisse) sur~$S$ en $x$, et
que le morphisme $f_{s}\colon X_{s}\to Y_{s}$ induit sur les fibres de
$s$ soit lisse en $x$.
\end{corollaire}

Seule la suffisance demande une d\'emonstration, et r\'esulte du
crit\`ere~\ref{II.2.1}, joint au crit\`ere de
platitude~(I~\ref{I.5.9}).

Pour
\'enoncer le r\'esultat suivant, \og rappelons\fg qu'un morphisme
$f\colon X\to Y$ localement de type fini est dit
\emph{équidimensionnel}
\index{equidimensionnel@\'equidimensionnel (morphisme)|hyperpage}en le point $x\in X$ si (posant $y=f(x)$) on peut trouver un voisinage
ouvert $U$ de~$x$, dont toute composante domine une composante de~$Y$
tel que, pour tout $y'\in Y$, les composantes irr\'eductibles de
$f^{-1}(y')\cap U$ aient toutes une m\^eme dimension
ind\'ependante de~$y'$. Il suffit d'ailleurs dans cette condition de
prendre pour $y'$ les points g\'en\'eriques des composantes
irr\'eductibles de~$Y$ passant par $y$, \emph{et} le point $y$. Si
par exemple $X$ et $Y$ sont int\`egres et $f$ dominant, la condition
signifie que les composantes des $f^{-1}(y)$ passant par $x$ ont \og la
bonne\fg dimension, \ie la dimension de la fibre g\'en\'erique
(rappelons qu'elles sont toujours $\ge$ la dimension de la fibre
g\'en\'erique). Si $f$ est \'equidimensionnel en $x$, la
dimension de sa fibre en $x$ \'etant $n$, et si $g\colon U\to
Y'=Y[t_{1},\dots,t_{n}]$ est un $Y$-morphisme d'un voisinage~$U$ de
$x$, induisant un morphisme sur les fibres de~$y$ qui est quasi-fini
en $x$ (ou encore, ce qui revient au m\^eme, si $g$ est quasi-fini
en~$x$), alors on montre que toute composante irr\'eductible de~$U$
passant par $x$ domine une composante irr\'eductible de~$Y'$. D'ailleurs en vertu du \og lemme de normalisation\fg, un tel $g$
existe toujours (et r\'eciproquement, s'il existe un $Y$-morphisme
quasi-fini $g$ d'un voisinage ouvert $U$ de~$x$ dans un $Y$-sch\'ema
de la forme $Y'=Y[t_{1},\dots,t_{n}]$, tel que toute composante de~$U$ passant par $x$ domine une composante de~$Y'$, alors $f$ est
\'equidimensionnel en~$x$). Ceci pos\'e:

\begin{proposition}
\label{II.2.3}
Soient $f\colon X\to Y$ un morphisme localement de type fini, $x$ un
point de~$X$, $y=f(x)$, on suppose $\cal{O}_{y}$ normal. Pour que $f$
soit lisse en $x$, il faut et il suffit que $f$ soit
\'equidimensionnel en $x$, et que $f^{-1}(y)$ soit lisse sur~$\kres(y)$
en $x$.
\end{proposition}

On voit aussit\^ot sur la d\'efinition qu'un morphisme lisse est
\'equidimensionnel (N.B. un morphisme plat de type fini n'est pas
n\'ecessairement \'equidimensionnel en $x$, m\^eme si sa fibre
en $x$ est irr\'eductible). Prouvons la r\'eciproque. Comme
$f^{-1}(y)$ est lisse sur~$\kres(y)$ en $x$, on peut supposer (rempla\c cant au besoin $X$ par un voisinage convenable de~$x$) qu'il existe
un $Y$-morphisme
$$
g\colon X\to Y[t_{1},\dots,t_{n}]=Y'
$$
induisant un morphisme \'etale sur les fibres de~$y$, et a fortiori
quasi-fini en $x$. Donc
$g$ est non ramifi\'e, et ($f$ \'etant \'equidimensionnel en
$x$) les composantes irr\'eductibles de~$X$ passant par $x$ dominent
chacun une composante de~$Y'$, a fortiori l'homomorphisme
$\cal{O}_{y'}\to\cal{O}_{x}$ d\'eduit de~$g$ (o\`u $y'=g(x)$) est
\emph{injectif}. Cet homomorphisme est de plus non ramifi\'e, et
$\cal{O}_{y'}$ est normal puisque localis\'e de l'anneau
$\cal{O}_{y}[t_{1},\dots,t_{n}]$, qui est normal puisque
$\cal{O}_{y}$ l'est. Donc l'homomorphisme $\cal{O}_{y'}\to\cal{O}_{x}$
est \'etale (I~\ref{I.9.5}~(ii)).

\begin{remarques}
\label{II.2.4}
L'\'enonc\'e pr\'ec\'edent vaut encore en rempla\c cant
l'hypoth\`ese que $\cal{O}_{y}$ est normal par l'hypoth\`ese plus
faible que $Y$ est \emph{géométriquement unibranche} en $y$,
(\cf I~\ref{I.11}) --- puisque (I~\ref{I.9.5}) vaut sous cette
hypoth\`ese. Profitons de l'occasion pour signaler en m\^eme temps
que si le corps r\'esiduel d'un anneau local int\`egre $A$ est
alg\'ebriquement clos, alors analytiquement int\`egre
\index{analytiquement int\`egre|hyperpage}(\ie $\widehat{A}$ est int\`egre) implique g\'eom\'etriquement
unibranche, la r\'eciproque \'etant vraie de plus dans toute
cat\'egorie de \og bons anneaux\fg, de fa\c con pr\'ecise dans une
cat\'egorie d'anneaux stable par les op\'erations usuelles, et
o\`u la compl\'etion d'un anneau local normal est normale
(condition remplie, en vertu du \og th\'eor\`eme de normalit\'e
analytique\fg de Zariski, dans la cat\'egorie des alg\`ebres
affines et leurs localis\'ees)\footnote{\Cf EGA IV~7.8.}.

\og Rappelons\fg enfin dans le contexte actuel le r\'esultat suivant,
d\^u \`a Hironaka\footnote{\Cf EGA IV 5.12.10.} qui permet parfois
de s'assurer que $f^{-1}(y)$ est un sch\'ema r\'eduit, \ie que
c'est aussi ce que de nombreux g\'eom\`etres alg\'ebristes
consid\'eraient abusivement comme la fibre (sans multiplicit\'e)
de~$f$ au-dessus de~$x$ (savoir $f^{-1}(y)_{\red}$):
\end{remarques}

\begin{proposition}
\label{II.2.5}
Soient $f\colon X\to Y$ un morphisme dominant de type fini de
pr\'esch\'emas r\'eduits, $y$ un point de~$Y$ tel que
$\cal{O}_{y}$ soit r\'egulier. On suppose que toutes les composantes
de~$f^{-1}(y)$ sont de multiplicit\'e~$1$
\index{composante de multiplicit\'e~$1$|hyperpage}(\cf d\'efinition plus bas), et que $f^{-1}(y)_{\red}$ est
normal. Alors $f^{-1}(y)$ est r\'eduit donc normal, $X$ est normal
en tous les points de~$f^{-1}(y)$, enfin $X$ est plat sur~$Y$ en tous
les points de~$f^{-1}(y)$.
\end{proposition}
On
dit qu'une composante $Z$ de~$f^{-1}(y)$ est de
\emph{multiplicité}~$1$
\index{multiplicit\'e d'une composante|hyperpage}si, $x$ d\'esignant le point g\'en\'erique de~$Z$, on a (i)
$\dim\cal{O}_{x}=\dim\cal{O}_{y}$ (\ie $Z$ n'est pas \og composante
exc\'edentaire\fg, c'est-\`a-dire n'est pas \og de dimension trop
grande\fg; (ii) l'id\'eal maximal de~$\cal{O}_{x}$ est engendr\'e
par l'id\'eal maximal de~$\cal{O}_{y}$ (qui a priori, en vertu du
choix de~$x$, engendre un id\'eal de d\'efinition
de~$\cal{O}_{x}$).

Compte tenu de~\ref{II.2.3} ou de~\ref{II.2.1} on trouve donc:

\begin{corollaire}
\label{II.2.6}
Soient $f\colon X\to Y$ un morphisme dominant de type fini de
pr\'esch\'emas r\'eduits, $y$ un point de~$Y$ tel que
$\cal{O}_{y}$ soit r\'egulier. Pour que $f$ soit lisse aux points de
$X$ au-dessus de~$y$, il faut et il suffit que les composantes de
$f^{-1}(y)$ soient de multiplicit\'es $1$, et que $f^{-1}(y)_{\red}$
soit lisse sur~$\kres(y)$.
\end{corollaire}

Cette situation \'etait surtout consid\'er\'ee par le pass\'e
quand $Y$ \'etait le spectre d'un anneau de valuation discr\`ete
$A$, et \'etait d\'esign\'ee commun\'ement sous des vocables
tels que: \og si la r\'eduction de~$X$ par rapport \`a la valuation
donn\'ee est jolie\fg... De plus, $X$ d\'esignait alors un
sous-sch\'ema (si on peut dire) ferm\'e d'un $\PP_{K}^{n}$ ($K$
\'etant le corps des fractions de~$A$) et faute d'un
langage
ad\'equat, le r\^ole plus intrins\`eque d'un objet \og d\'efini
sur~$A$\fg (et non seulement sur~$K$) n'apparaissait gu\`ere.

\section{Propri\'et\'es de permanence}
\label{II.3}

\begin{proposition}
\label{II.3.1}
Soit $f\colon X\to Y$ un morphisme, soit $x\in X$ et
$y=f(x)$. Supposons $f$ lisse en~$x$. Pour que $\cal{O}_{x}$ soit
r\'eduit (\resp r\'egulier, \resp normal) il faut et il suffit que
$\cal{O}_{y}$ le soit.
\end{proposition}

Cet \'enonc\'e est en effet connu quand $X$ est de la forme
$Y[t_{1},\dots,t_{n}]$, et il est d\'emontr\'e dans (I, \no \ref{I.9})
pour un morphisme \'etale; le cas g\'en\'eral s'en d\'eduit
aussit\^ot gr\^ace \`a la d\'efinition~\ref{II.1.1}.

Nous ne d\'etaillons pas ici les autres propri\'et\'es de
permanence, r\'esultant d\'ej\`a de la seule platitude, ou du
fait que $X$ est localement quasi-fini et plat au-dessus d'un
$Y$-pr\'esch\'ema de la forme $Y[t_{1},\dots,t_{n}]$ (ou, comme
nous dirons, que $X$ est
Cohen-Macaulay
\index{Cohen-Macaulay (sch\'ema de)|hyperpage}au-dessus de~$Y$). Signalons seulement que de ce dernier fait
r\'esulte que
\begin{equation}
\label{eq:II.3.1}
{\dim\cal{O}_{x}=\dim\cal{O}_{y}+n-d,\quad
\prof\cal{O}_{x}=\prof\cal{O}_{y}+n-d},
\end{equation}
o\`u $n$ est la dimension de la fibre de~$f$ en $x$, et $d$ le
degr\'e de transcendance de~$\kres(x)$ sur~$\kres(y)$, d'o\`u (posant
$\coprof=\dim-\prof$)\footnote{Pour ces formules, \cf EGA~IV 6.1
et~6.3.}
\begin{equation}
\label{eq:II.3.2}
{\coprof\cal{O}_{x}=\coprof\cal{O}_{y}}.
\end{equation}
Il en r\'esulte par exemple que $\cal{O}_{x}$ est Cohen-Macaulay
(\resp sans composantes immerg\'ees) si et seulement si il en est de
m\^eme de~$\cal{O}_{y}$.

\section{Propri\'et\'es diff\'erentielles des morphismes lisses}
\label{II.4}
Pour simplifier, nous nous restreindrons pour l'essentiel au calcul
diff\'erentiel d'ordre~$1$, nous bornant \`a de rapides
indications pour l'ordre sup\'erieur (o\`u les r\'esultats sont
tout aussi simples).

Pour la d\'efinition du faisceau des $1$-diff\'erentielles
$\it{\Omega}^1_{X/Y}$ d'un $Y$-pr\'esch\'ema $X$,
\cf (I~\No\ref{I.1}). Supposons que $X$ et $Y$ soient des
$S$-pr\'esch\'emas, le morphisme structural $f\colon X\to Y$
\'etant un $S$-morphisme. Alors $f$ d\'efinit un homomorphisme de
Modules (compatible avec $f$)
\begin{equation}
\label{eq:II.4.1}
{f^*\colon \it{\Omega}^1_{Y/S}\to\it{\Omega}^1_{X/S}}
\end{equation}
en d'autres termes, $\it{\Omega}^1_{X/S}$ est \emph{contravariant} en
le $S$-pr\'esch\'ema $X$. D'ailleurs~\eqref{eq:II.4.1}
\'equivaut \`a un homomorphisme de Modules sur~$X$
\begin{equation*}
\label{eq:II.4.1bis}
\tag{4.1\textrm{bis}}
{f^*\big(\it{\Omega}^1_{Y/S}\big)\to\it{\Omega}^1_{X/S}}
\end{equation*}
\'egalement d\'enot\'e par $f^*$ \`a d\'efaut de mieux, et
qui s'ins\`ere dans une suite exacte canonique d'homomorphismes de
Modules
\begin{equation}
\label{eq:II.4.2}
{f^*\big(\it{\Omega}^1_{Y/S}\big)\to
\it{\Omega}^1_{X/S}\to\it{\Omega}^1_{X/Y}\to 0}
\end{equation}
Tous ces homomorphismes sont d\'efinis par la condition d'\^etre de
nature locale (ce qui ram\`ene au cas affine) et de commuter avec
les op\'erateurs $\rd$. L'exactitude de~\eqref{eq:II.4.2} est
classique et triviale, et se transcrit dans le cas affine en la suite
exacte (correspondant \`a un homomorphisme $B\to C$ de
$A$-alg\`ebres):
\begin{equation*}
\label{eq:II.4.2bis}
\tag{4.2\textrm{bis}}
{\Omega^1_{B/A}\otimes_{B}C\to\Omega^1_{C/A}\to\Omega^1_{C/B}\to 0}
\end{equation*}

\begin{lemme}
\label{II.4.1}
Soit
$f\colon X\to Y$ un morphisme de~$S$-pr\'esch\'emas. Si $f$
est non ramifi\'e (\resp \'etale) alors
$f^*\big(\it{\Omega}^1_{Y/S}\big)\to\it{\Omega}^1_{X/S}$
est surjectif (\resp un isomorphisme). La r\'eciproque est vraie
dans le cas \og non ramifi\'e\fg, si $f$ est suppos\'e localement de
type fini.
\end{lemme}

Le cas non ramifi\'e r\'esulte de la suite
exacte~\eqref{eq:II.4.2} et de (I~\ref{I.3.1}), mais peut aussi se
voir directement comme le cas \'etale. Consid\'erons le diagramme
\begin{equation*}
\xymatrix@C=1.5cm{{X}\ar[r]^-{\Delta_{X/Y}}&{X\times_{Y}X}\ar[d]\ar[r]&{X\times_{S}X}\ar[d]\\ &{Y}\ar[r]^-{\Delta_{Y/S}}&{Y\times_{S}Y}}
\end{equation*}
dans lequel $X\times_{Y}X$ s'identifie au produit fibr\'e de~$Y$ et
$X\times_{S}X$ sur~$Y\times_{S}Y$. Comme $f$ est non ramifi\'e,
$X\to X\times_{Y}X$ est une immersion ouverte, donc le faisceau
\og conormal\fg de l'immersion compos\'ee $\Delta_{X/S}$ de cette
derni\`ere avec $X\times_{Y}X\to X\times_{S}X$ est isomorphe \`a
l'image inverse sur~$X$ du faisceau conormal pour l'immersion
$X\times_{Y}X\to X\times_{S}X$. D'autre part, $X\to Y$ \'etant
\'etale donc plat, $X\times_{S}X\to Y\times_{S}Y$ est plat, donc le
faisceau conormal pour l'immersion $X\times_{Y}X\to X\times_{S}X$ est
isomorphe \`a l'image inverse du faisceau conormal pour l'immersion
$Y\to Y\times_{S}Y$, \ie l'image inverse de~$\it{\Omega}^1_{Y/S}$. La
conclusion en r\'esulte.

\begin{lemme}
\label{II.4.2}
Soit $X=Y[t_{1},\dots,t_{n}]$, $Y$ \'etant un
$S$-pr\'esch\'ema. Alors la suite d'homomorphismes canoniques
$$
0\to f^*\big(\it{\Omega}^1_{Y/S}\big)\to
\it{\Omega}^1_{X/S}\to\it{\Omega}^1_{X/Y}\to 0
$$
est exacte et $\it{\Omega}^1_{X/Y}$ est libre de base les
$\rd_{X/Y}t_{i}$.
\end{lemme}

La v\'erification (purement affine) est imm\'ediate. (N.B. on
conna\^it d\'ej\`a l'exactitude de~\eqref{eq:II.4.2}).

Combinant ces deux \'enonc\'es et d\'efinition~\ref{II.1.1}, on
trouve

\begin{theoreme}
\label{II.4.3}
Soient $f\colon X\to Y$ un morphisme lisse de~$S$-pr\'esch\'emas,
alors:
\begin{enumerate}
\item[(i)] La suite d'homomorphismes canoniques
$$
0\to f^*\big(\it{\Omega}^1_{Y/S}\big)\to
\it{\Omega}^1_{X/S}\to\it{\Omega}^1_{X/Y}\to 0
$$
est exacte.
\item[(ii)] $\it{\Omega}^1_{X/Y}$ est localement libre, son rang $n$
en $x$ est \'egal \`a la dimension relative de~$f$ en~$x$.
\end{enumerate}
\end{theoreme}

\begin{corollaire}
\label{II.4.4}
L'homomorphisme
$f^*\big(\it{\Omega}^1_{Y/S}\big)\to\it{\Omega}^1_{X/S}$ est
injectif, son image dans $\it{\Omega}^1_{X/S}$ est localement facteur
direct.
\end{corollaire}

Soit $u\colon F\to G$ un homomorphisme de Modules sur le
pr\'esch\'ema $X$, on dit qu'il est \emph{universellement
injectif}
\index{universellement injectif (morphisme de modules)|hyperpage}en $x\in X$, si l'homomorphisme $F_{x}\to G_{x}$ de
$\cal{O}_{x}$-modules est injectif, et reste tel par tensorisation
avec toute $\cal{O}_{x}$-alg\`ebre (ou, ce qui revient au m\^eme
d'ailleurs, avec tout $\cal{O}_{x}$-module). Il suffit par exemple
qu'il existe un voisinage ouvert $U$ de~$x$ tel que $u$ induise un
isomorphisme de~$F|U$ sur un facteur direct de~$G|U$, cette condition
est aussi n\'ecessaire lorsque $F$ et $G$ sont libres (et de type
fini) dans un voisinage de~$x$, de fa\c con pr\'ecise dans ce cas
les conditions suivantes sont \'equivalentes:
\begin{enumerate}
\item[(i)] $u$ est injectif en $x$ et $\Coker u$ libre en $x$;
\item[(ii)] Il existe un voisinage ouvert $U$ de~$x$ tel que $u$
induise un isomorphisme de~$F|U$ sur un facteur direct de~$G|U$;
\item[(iii)] $u$ est universellement injectif en $x$;
\item[(iv)] l'homomorphisme
$F_{x}\otimes \kres(x)\to G_{x}\otimes \kres(x)$
sur les \og fibres\fg restreintes induit par $u$ est injectif;
\item[(v)] L'homomorphisme transpos\'e $\check{G}\to\check{F}$ est
surjectif au point $x$ (ou encore, ce qui revient au m\^eme, au
voisinage de~$x$).
\end{enumerate}
(D\'emonstration circulaire, (iv)$\To$(v) r\'esulte de
Nakayama, d'autre part (v)$\To$(i)
puisqu'un
faisceau
quotient localement libre est n\'ecessairement facteur
direct). G\'eom\'etriquement, la situation envisag\'ee signifie
que $u$ correspond \`a un isomorphisme du fibr\'e vectoriel dont
le faisceau des sections
est~$F$,
sur un sous-fibr\'e du
fibr\'e vectoriel analogue d\'efini par $G$. Bien entendu, il ne
suffit pas pour cela que
$F\to G$
soit injectif.

\begin{corollaire}
\label{II.4.5}
Soit $f\colon X\to Y$ un morphisme de~$S$-pr\'esch\'emas,
localement de type fini, $x\in X$, $y=f(x)$, $s$ la projection de~$x$
et $y$ sur~$S$. On suppose $Y$ lisse en $y$ sur~$S$. Conditions
\'equivalentes:
\begin{enumerate}
\item[(i)] $f$ est lisse en $x$.
\item[(ii)] $X$ est lisse sur
$S$
en~$x$, et
$f^*\big(\it{\Omega}^1_{Y/S}\big)\to\it{\Omega}^1_{X/S}$ est
universellement injectif en~$x$, \ie c'est un homomorphisme injectif en~$x$ et son conoyau $\it{\Omega}^1_{X/Y}$ est libre
en~$x$.
\end{enumerate}
\end{corollaire}

La n\'ecessit\'e r\'esulte de~\ref{II.1.3} (iii) et
de~\ref{II.4.3} (i)\,(ii), prouvons la suffisance. Comme les
$\rd g$ ($g\in\cal{O}_{x}$) engendrent le module
$\it{\Omega}^1_{X/Y}$ en $x$, on peut trouver des
$g_{i}$ ($1\le i\le n$) tels que les images des $\rd g_{i}$ dans
$\big(\it{\Omega}^1_{X/Y}\big)_{x}$ forment une base de ce
module. Prenant $X$ assez petit, on peut supposer que les $g_{i}$
proviennent de sections de~$\cal{O}_{X}$, et d\'efinissent donc un
$Y$-morphisme $g\colon X\to Y'=Y[t_{1},\dots,t_{n}]$. Utilisant
l'hypoth\`ese et lemme~\ref{II.4.2}, on voit facilement que
l'homomorphisme correspondant
$g^*\big(\it{\Omega}^1_{Y'/S}\big)\to\it{\Omega}^1_{X/S}$ est
bijectif en $x$, ce qui nous ram\`ene \`a prouver le

\begin{corollaire}
\label{II.4.6}
Soit $f\colon X\to Y$ un morphisme de~$S$-pr\'esch\'emas
lisses. Pour que $f$ soit \'etale en $x\in X$, il faut et il suffit
que $f^*\big(\it{\Omega}^1_{Y/S}\big)\to\it{\Omega}^1_{X/S}$ soit
un isomorphisme en $x$.
\end{corollaire}

On sait que c'est n\'ecessaire par~\ref{II.4.1}, et cette condition
implique que $f$ est non ramifi\'e en $x$ par le m\^eme lemme. En
vertu de~\ref{II.2.2}, on est ramen\'e au cas o\`u
$S=\Spec(k)$. Comme $Y$ est lisse sur~$k$, il est r\'egulier, donc a
fortiori normal, et en vertu de (I~\ref{I.9.5} (ii)) on est ramen\'e
\`a prouver que $\cal{O}_{y}\to\cal{O}_{x}$ est injectif, ou encore
que $\cal{O}_{y}$ et $\cal{O}_{x}$ ont m\^eme dimension. Or ces
dimensions sont respectivement les rangs de~$\it{\Omega}^1_{Y/k}$ et
$\it{\Omega}^1_{X/k}$ en $y$ \resp $x$, donc \'egaux en vertu de
l'hypoth\`ese.
\begin{remarques}
\label{II.4.7}
$X$ et $Y$ \'etant suppos\'es lisses sur~$S$, le
crit\`ere~\ref{II.4.5} (ii) de lissit\'e de~$f\colon X\to Y$ peut
encore s'\'enoncer en disant que pour tout $x\in X$, l'application
\emph{tangente} (relativement \`a la base $S$) de~$f$ en $x$,
\ie la transpos\'ee de l'homomorphisme des $\kres(x)$ espaces vectoriels
de dimension finie, fibres restreintes de
$f^*\big(\it{\Omega}^1_{Y/S}\big)$ et $\it{\Omega}^1_{X/S}$ en
$x$, est \emph{surjective}. C'est l\`a une hypoth\`ese bien
famili\`ere en particulier parmi les gens travaillant avec les
espaces analytiques. L'hypoth\`ese de non singularit\'e qu'ils
font d'ordinaire (qui signifie que $X$ et $Y$ sont \og lisses sur~$\CC$\fg, \cf \No \ref{II.5}) ne semble due qu'\`a la peur qu'inspirent
encore \`a bien des g\'eom\`etres les points singuliers des
vari\'et\'es alg\'ebriques ou espaces analytiques.
\end{remarques}

Signalons le cas particulier suivant de~\ref{II.4.6}:
\begin{corollaire}
\label{II.4.8}
Soient $X$ un $S$-pr\'esch\'ema, $g\colon X\to
S[t_{1},\dots,t_{n}]$ un $S$-morphisme, d\'efini par les sections
$g_{i}$ ($1\le i\le n$) de~$\cal{O}_{X}$, $x$ un point de~$X$ tel que
$X$ soit lisse sur~$S$ en $x$. Pour que $g$ soit \'etale en $x$, il
faut et il suffit que les $\rd g_{i}$
($1\le i\le n$) forment une base de~$\it{\Omega}^1_{X/S}$ en $x$ (ou,
ce qui revient au m\^eme, que leurs images dans
$\it{\Omega}^1_{X/S}(x) =
\big(\it{\Omega}^1_{X/S}\big)_{x}\otimes_{\cal{O}_{x}}\kres(x)$ forment
une base de cet espace vectoriel sur~$\kres(x)$).
\end{corollaire}

Soient $X$ un pr\'esch\'ema, $Y$ un sous-pr\'esch\'ema
ferm\'e de~$X$ d\'efini par un faisceau coh\'erent $\cal{J}$
d'id\'eaux. Donc $\cal{J}/\cal{J}^2$ peut \^etre consid\'er\'e
comme un faisceau coh\'erent sur~$Y$ (le \emph{faisceau conormal}
\index{conormal (faisceau)|hyperpage}\index{faisceau conormal|hyperpage}de~$Y$ dans~$X$). Si maintenant $X$ est un $S$-pr\'esch\'ema, on a
une suite exacte canonique de faisceaux quasi-coh\'erents sur~$Y$
\begin{equation}
\label{eq:II.4.3}
{\cal{J}/\cal{J}^2\lto{\rd}
\it{\Omega}^1_{X/S}\otimes_{\cal{O}_{X}}\cal{O}_{Y}\to
\it{\Omega}^1_{Y/S}\to 0}
\end{equation}
dont la partie de droite n'est autre que~\eqref{eq:II.4.2} (avec le
r\^ole de~$X$ et $Y$ interchang\'es, compte tenu que
$\it{\Omega}^1_{Y/X}=0$), tandis que l'homomorphisme
$\cal{J}/\cal{J}^2\to\it{\Omega}^1_{X/S}\otimes_{\cal{O}_{X}}\cal{O}_{Y}$
est d\'eduit de l'homomorphisme (en g\'en\'eral non
lin\'eaire) $g\to\rd g$ par passage aux
quotients. L'exactitude de~\eqref{eq:II.4.3} est classique et
d'ailleurs triviale, et s'interpr\`ete dans le cas affine par la
suite exacte suivante (correspondante \`a un homomorphisme surjectif
$B\to C$ de~$A$-alg\`ebres, de noyau $J$):
\begin{equation*}
\label{eq:II.4.3bis}
\tag{4.3\textrm{bis}}
{J/J^2\to\Omega^1_{B/A}\otimes_{B}C\to\Omega^1_{C/A}\to 0\qquad
(C=B/J)}
\end{equation*}
(suite exacte qui avait d\'ej\`a \'et\'e utilis\'ee
implicitement dans la d\'emonstration de (I~\ref{I.7.2})!).

\begin{proposition}
\label{II.4.9}
Soient $X$ un $S$-pr\'esch\'ema, $Y$ un sous-pr\'esch\'ema
ferm\'e de~$X$ d\'efini par un faisceau coh\'erent $\cal{J}$
d'id\'eaux sur~$X$, $x$ un point de~$X$, $g_i$ ($1 \leq i \leq n$)
des sections de~$\cal{O}_X$, d\'efinissant un $S$-morphisme
$$
g \colon X \to S[t_1,\dots,t_n] = X'
$$
enfin $p$ un entier, $0 \leq p \leq n$. On suppose $X$ \emph{lisse sur~$S$ en $x$}. Les conditions suivantes sont \'equivalentes:
\begin{enumerate}
\item[(i)] Il existe un voisinage ouvert $X_1$ de~$x$ dans $X$ tel que
$g|X_1$ soit \emph{étale} et que $Y_1=Y \cap X_1$ (trace de~$Y$
sur~$X_1$) soit \emph{l'image inverse} du sous-pr\'esch\'ema
ferm\'e $Y'=S[t_{p+1},\dots,t_n]$
de
$X'=S[t_1,\dots,t_n]$ (\ie les $g_i$ ($1 \leq i \leq p$) engendrent
$\cal{J}|X_1$):
$$
\xymatrix{{Y_1}\ar[d]\ar[r]&{X_1}\ar[d]^{\'etale}\\{Y'=S[t_{p+1},\dots, t_n]}\ar[r]&{X'=S[t_1,\dots,t_n]}}
$$
\item[(ii)]
$Y$
est \emph{lisse sur~$S$ en $x$}, les $g_i$ ($1 \leq i \leq p$)
d\'efinissent des \emph{éléments de} $\cal{J}_x$, les $d
g_i(x)$ ($1 \leq i \leq n$) forment une \emph{base de}
$\it{\Omega}^1_{X/S}(x)$ sur~$\kres(x)$, les $dg'_i(x)$ ($p+1 \leq i \leq
n$) forment une \emph{base de} $\it{\Omega}^1_{Y/S}(x)$ sur~$\kres(x)$
(o\`u les $g'_i$ d\'esignent les restrictions des $g_i$ \`a~$Y$;
les diff\'erentielles sont prises par rapport \`a~$S$).
\item[(iii)] Les $g_i$ ($1 \leq i \leq p$) d\'efinissent un
\emph{système de générateurs} de~$\cal{J}_x$, et les
$dg_i(x)$ ($1 \leq i \leq n$) forment une \emph{base de}
$\it{\Omega}^1_{X/S}(x)$ sur~$\kres(x)$.
\item[(iv)] $Y$ est \emph{ lisse sur~$S$} en $x$, les $g_i$
($1 \leq i \leq p$)
forment un \emph{système minimal de générateurs de}
$\cal{J}_x$, les $dg'_i(x)$ ($p+1 \leq i \leq n$) forment une
\emph{base de} $\it \Omega^1_{Y/S}(x)$ sur~$\kres(x)$.
\end{enumerate}

De plus, sous ces conditions, $\cal{J}/\cal{J}^2$ est un Module libre
sur~$Y$ en $x$, admettant comme \emph{base en $x$} les classes des
$g_i$ ($1 \leq i \leq p$), \emph{et l'homomorphisme} canonique
$\cal{J}/\cal{J}^2 \to \it{\Omega}^1_{X/S} \otimes \cal{O}_Y$ est
\emph{universellement injectif en $x$}.
\end{proposition}

\begin{remarquestar}
Cela implique que $p$ est bien d\'etermin\'e par les autres
conditions, soit comme \emph{rang} du Module libre
$\cal{J}/\cal{J}^2$ sur~$Y$ en $x$, ou encore le \emph{nombre minimum
de générateurs} de~$\cal{J}_x$ sur~$X$, ou enfin par le fait
que la dimension relative de~$Y$ rel. $S$ en $x$ est $n-p$.
\end{remarquestar}

\subsubsection*{D\'emonstration}
Supposons d'abord (i) v\'erifi\'e. Alors par (I~\ref{I.4.6}
(iii)) $Y_1$ est \'etale sur~$Y'$, donc par d\'efinition il est
lisse sur~$S$ en $x$ (de dimension relative $n-p$), il en est donc de
m\^eme de~$Y$. Il r\'esulte alors de \eqref{II.4.8} que les $dg_i$
($1 \leq i \leq n$) forment une base de~$\it{\Omega}^1_{X/S}$ en $x$,
et que les $dg'_i$ ($p+1 \leq i \leq n$) une base de
$\it{\Omega}^1_{Y/S}$ en $x$, d'o\`u il r\'esulte par la suite
exacte \eqref{eq:II.4.3} que les $g_i$ ($1 \leq i \leq p$) sont
lin\'eairement ind\'ependants dans $\cal{J}/\cal{J}^2$
(consid\'er\'e comme Module sur~$Y$) en $x$; comme les $g_i$ ($1
\leq i \leq p$) engendrent $\cal{J}_x$, il s'ensuit que les $g_i \mod
\cal{J}_x^2$ forment une \emph{base de} $\cal{J}/\cal{J}^2$ en
$x$. Cela implique d'une part que les $g_i$ ($1 \leq i \leq p$)
forment un syst\`eme \emph{minimal} de g\'en\'erateurs de
$\cal{J}_x$, d'autre part que l'homomorphisme $\cal{J}/\cal{J}^2 \to
\it{\Omega}^1_{X/S} \otimes \cal{O}_Y$ de \eqref{eq:II.4.3} est
universellement injectif en $x$ (car applique une base d'un Module
libre en $x$ sur une partie d'une base d'un Module libre en $x$ --- N.B. il s'agit de~$Y$-Modules). Cela prouve que (i) implique (ii),
(iii), (iv), ainsi que les derni\`eres assertions de
proposition~\ref{I.4.9}.

(iii) implique (i) en vertu de corollaire~\ref{I.4.8}.

(ii)
implique (i). En effet, la premi\`ere hypoth\`ese dans (ii)
signifie que (quitte \`a remplacer $X$ par un voisinage ouvert de
$x$ dans $X$) $g$ induit un morphisme $h\colon Y \to
Y'$. D'apr\`es~\ref{I.4.8}, les deux autres hypoth\`eses de (ii)
signifient que $g$ est \'etale en $x$, et $h$ \'etale en $x$. Soit
alors $Y''$ l'image inverse de~$Y'$ par $g$. Donc $Y$ est un
sous-pr\'esch\'ema ferm\'e de~$Y''$, qui est \'etale sur~$Y'$
en $x$ par (I~\ref{I.4.6}~(iii)) puisque $g$ est \'etale en
$x$. Donc le morphisme d'immersion $Y \to Y''$ est lui-m\^eme
\'etale (I~\ref{I.4.8}) donc une immersion ouverte (I~\ref{I.5.8} ou
I~\ref{I.5.2}), donc rempla\c cant encore $X$ par un voisinage
ouvert convenable $X_1$ de~$x$, on obtient (i).

Ce qui pr\'ec\`ede \'etablit l'\'equivalence des conditions
(i), (ii), (iii), et le fait qu'elles impliquent (iv), il reste \`a
prouver que (iv)$\To$(ii), ce qui est imm\'ediat (compte
tenu que $\it{\Omega}^1_{X/S}$ est libre sur~$X$ en $x$) une fois
qu'on sait que le fait que $Y$ est lisse sur~$S$ en $x$ implique que
$\cal{J}/\cal{J}^2$ est libre sur~$Y$ en $x$, et l'homomorphisme
$\cal{J}/\cal{J}^2 \to \it{\Omega}^1_{X/S}\otimes \cal{O}_Y$
universellement injectif en $x$. Ce dernier point est inclus dans le
\begin{theoreme}
\label{II.4.10}
Soient $X$ un $S$-pr\'esch\'ema lisse, $Y$ un
sous-pr\'esch\'ema ferm\'e de~$X$ d\'efini par un faisceau
coh\'erent $\cal{J}$ d'id\'eaux sur~$X$, $x$ un point de~$X$. Les
conditions suivantes sont \'equivalentes:
\begin{enumerate}
\item[(i)] $Y$ est \emph{lisse sur~$S$ en $x$}.
\item[(ii)] Il existe un voisinage ouvert $X_1$ de~$x$ dans $X$ et un
$S$-morphisme \emph{étale}
$$
g\colon X_1 \to X'=S[t_1,\dots,t_n]
$$
tel que $Y_1=Y \cap X_1$ (trace de~$Y$ sur~$X_1$) soit le
sous-pr\'esch\'ema de~$X_1$ \emph{image inverse} par $g$ du
sous-pr\'esch\'ema ferm\'e $Y'=S[t_{p+1},\dots,t_n]$ de
$X'=S[t_1,\dots,t_n]$, pour un $p$ convenable.
\item[(iii)] Il existe des \emph{générateurs $g_i$ ($1 \leq i
\leq p$) de~$\cal{J}_x$}, tels que les $dg_i$ forment une partie d'une
base de~$\it{\Omega}^1_{X/S}$ en $x$ (ou, ce qui revient au m\^eme,
tel que les $dg_i(x)$ dans $\it{\Omega}^1_{X/S}$ soient
lin\'eairement ind\'ependant sur~$\kres(x)$).
\item[(iv)] Le faisceau $\cal{J}/\cal{J}^2$ est libre sur~$Y$ en $x$,
et l'homomorphisme canonique
$$
d\colon \cal{J}/\cal{J}^2 \to \it{\Omega}^1_{X/S}\otimes \cal{O}_Y
$$
est universellement injectif en~$x$; ou encore: la suite
d'homomorphismes canoniques
$$
0 \to \cal{J}/\cal{J}^2 \to \it{\Omega}^1_{X/S}\otimes \cal{O}_Y \to
\it{\Omega}^1_{Y/S} \to 0
$$
est exacte en $x$, et $\it{\Omega}^1_{Y/S}$ est localement libre
en~$x$.
\end{enumerate}
\end{theoreme}

\subsubsection*{D\'emonstration}
On
sait d\'ej\`a que (ii) implique (i), (iii), (iv) d'apr\`es ce qui
pr\'ec\`ede. Prouvons que (i)$\To$(ii) (ce qui
ach\`evera en m\^eme temps la d\'emonstration
de~\ref{I.4.9}).
En vertu de th\'eor\`eme \ref{II.4.3}~(ii), les deux derniers
termes dans la suite exacte \eqref{eq:II.4.3} sont des Modules libres
sur~$Y$. Donc, comme les images dans
$\it{\Omega}^1_{X/S}\otimes_{\cal{O}_X} \cal{O}_Y$ des $dg$ ($g \in
\cal{O}_X$) engendrent ce module en $x$, donc leurs images dans
$\it{\Omega}^1_{Y/S}$ engendrent ce dernier en~$x$, on peut trouver
des $g_i$ ($p+1 \leq i \leq n$) dans $\cal{O}_X$ tels que les $dg'_i$
forment une base de~$\it{\Omega}^1_{Y/S}$, puis (en vertu de
l'exactitude de \eqref{eq:II.4.3}) compl\'eter le syst\`eme des
$dg_i$ ($p+1 \leq i \leq n$) en une base du terme m\'edian par des
\'el\'ements de la forme $dg_i$ ($1 \leq i \leq n$) o\`u les
$g_i$ ($1 \leq i \leq p$) \emph{sont dans $\cal{J}_x$.} Les $g_i$
proviennent de sections de~$\cal{O}_X$ sur un voisinage de~$x$ dans
$X$, qu'on peut supposer \'egal \`a $X$. On est alors sous les
conditions de~\ref{II.4.8}~(ii), et on a \'etabli que cela implique
la condition~\ref{II.4.8}~(i), d'o\`u~\ref{II.4.10}~(ii).

L'implication (iii)$\To$(ii) dans~\ref{II.4.10} r\'esulte
aussit\^ot de l'implication (iii)$\To$(i)
dans~\ref{II.4.8}. Donc
(i), (ii), (iii)
sont \'equivalents, et
impliquent (iv). Enfin, il est trivial que (iv) implique (iii), compte
tenu que des $g_i \in \cal{J}_x$ qui forment une base de~$\cal{J}_x
\mod \cal{J}_x^2$ engendrent $\cal{J}_x$ (Nakayama).

De plus, la d\'emonstration qui pr\'ec\`ede montre ceci:
\begin{corollaire}
\label{II.4.11}
Soient $X$ un $S$-pr\'esch\'ema, $Y$ un sous-pr\'esch\'ema
ferm\'e d\'efini par un faisceau coh\'erent $\cal{J}$
d'id\'eaux sur~$X$, $x$ un point de~$Y$. On suppose \emph{$X$ et $Y$
lisses sur~$S$ en $x$}. Soient $g_i$ des sections de~$\cal{J}$ ($1
\leq i \leq p$). Les conditions suivantes sont \'equivalentes:
\begin{enumerate}
\item[(i)] Les $g_i$ \emph{engendrent} $\cal{J}_x$ et les $dg_i(x)$
sont \emph{linéairement indépendants} dans
$\it{\Omega}^1_{X/S}(x)$
sur~$\kres(x)$.
\item[(ii)] Les $g_i \mod \cal{J}^2$ forment une base de
$\cal{J}/\cal{J}^2$ en $x$.
\item[(iii)] Les $g_i$ forment un syst\`eme minimal de
g\'en\'erateurs de~$\cal{J}_x$.
\item[(iv)] On peut trouver d'autres sections $g_i$ ($p+1 \leq i \leq
n$) de~$\cal{O}_X$ sur un voisinage $X_1$ \emph{de} $X$,
d\'efinissant avec les pr\'ec\'edents un morphisme
\emph{étale} $X_1 \to X'=S[t_1,\dots,t_n]$ tel que $Y_1=Y \cap
X_1$ soit \emph{l'image inverse} par $g$ du sous-pr\'esch\'ema
ferm\'e $Y'=S[t_{p+1},\dots,t_n]$ de~$X'=S[t_{1},\dots,t_n]$.
\end{enumerate}
\end{corollaire}

En
particulier:

\begin{corollaire}
\label{II.4.12}
Soient $X$ un $S$-pr\'esch\'ema, $F$ une section de~$\cal{O}_X$,
$Y$ le sous-pr\'esch\'ema des z\'eros de~$F$ (d\'efini par
l'Id\'eal coh\'erent $F\cdot\cal{O}_X$), $x$ un point de~$Y$. On
suppose $X$ simple sur~$S$ en $x$. Pour que $Y$ soit lisse sur~$S$ en
$x$, il faut et il suffit que ou bien $F$ soit nul au voisinage de
$x$, ou bien que $dF(x) \neq 0$ (o\`u $dF(x)$ d\'esigne l'image de
$dF$ dans l'espace vectoriel $\it{\Omega}^1_{X/S}(x)$ sur~$\kres(x)$).
\end{corollaire}

C'est suffisant en vertu de~\ref{II.4.10} crit\`ere~(iii). C'est
n\'ecessaire, car comme $\cal{J}$ est engendr\'e par un
\'el\'ement, il faut d'abord que $\cal{J}/\cal{J}^2$ au point $x$
soit libre de rang $\leq 1$. Si ce rang est $0$,
\ie $\cal{J}/\cal{J}^2=0$ en $x$, il s'ensuit que $\cal{J}=0$ en $x$
par Nakayama, \ie $F$ est nul au voisinage de~$x$. Si ce rang est
$1$, alors $F$ forme un syst\`eme minimal de g\'en\'erateurs de
$\cal{J}$ en $x$, et on conclut par \eqref{II.4.11}, \'equivalence de
(i) et~(iii)).
\begin{corollaire}
\label{II.4.13}
Soient $Y$ un $S$-pr\'esch\'ema localement de type fini, $S'$ un
$S$-pr\'e\-sch\'ema \emph{plat}, $Y'=Y \times_S S'$, $x'$ un point
de~$Y'$ et $x$ son image canonique dans~$Y$. Pour que $Y$ soit lisse
sur~$S$ en~$X$, il faut et il suffit que $Y'$ soit lisse sur~$S'$ en~$x'$. En particulier, si $S'\to S$ est plat et surjectif, $Y$ est
lisse sur~$S$
si et seulement si\xspace $Y'$ est lisse sur~$S'$.
\end{corollaire}

Il n'y a \`a prouver que la suffisance (la n\'ecessit\'e a
\'et\'e vue dans~\ref{II.1.3}~(ii)). On peut supposer (rempla\c cant $Y$ par un voisinage convenable de~$x$, $Y'$ par l'image inverse
de ce dernier) que $Y$ est affine de type fini sur~$S$ affine, donc
$Y$ est isomorphe \`a un sous-pr\'esch\'ema ferm\'e d'un
sch\'ema $S[t_1,\dots,t_n]$. Par suite, $Y'$ s'identifie \`a un
sous-pr\'esch\'ema ferm\'e de~$X'=X \times_S S'$. Comme $X$ est
lisse sur~$S$, donc $X'$ lisse sur~$S'$, on peut appliquer les
crit\`eres de lissit\'e~\ref{II.4.10}. Ici, le crit\`ere (iv)
donne le r\'esultat aussit\^ot.

\begin{remarques}
\label{II.4.14}
Le crit\`ere (iii) de th\'eor\`eme~\ref{II.4.10} m\'erite
d'\^etre appel\'e \emph{critère jacobien de lissité}.
\index{crit\`ere jacobien de lissit\'e|hyperpage}\index{jacobien (crit\`ere de lissit\'e)|hyperpage}\index{lissit\'e (crit\`ere jacobien de)|hyperpage}Il permet de reconna\^{\i}tre, th\'eoriquement, si un
$S$-pr\'esch\'ema donn\'e $Y$ est lisse sur~$S$ en un point $x$
de~$Y$, puisque il existe toujours un voisinage de~$Y$ isomorphe \`a
un sous-pr\'esch\'ema d'un $S$-pr\'esch\'ema lisse $X$, par
exemple $X=S[t_1,\dots,t_n]$. C'est d'ailleurs pour
$X=S[t_1,\dots,t_n]$, $S=\Spec(A)$, qu'on \'enonce d'habitude le
crit\`ere jacobien (bien
entendu, dans le cas classique envisag\'e par Zariski, $A$ \'etait
un corps). On laisse au lecteur de donner l'\'enonc\'e relatif
\`a la donn\'ee d'un id\'eal $J$ de~$A[t_1,\dots,t_n]$ et d'un
id\'eal premier le contenant, auquel on est ainsi conduit. Notons
qu'il semble bien \`a l'heure actuelle (et surtout depuis que Nagata
est parvenu \`a g\'en\'eraliser par des m\'ethodes
non-diff\'erentielles le th\'eor\`eme de Zariski disant que
l'ensemble des points r\'eguliers d'un sch\'ema alg\'ebrique est
ouvert) que le crit\`ere jacobien n'a gu\`ere d'int\'er\^et
que sous la forme o\`u nous le donnons ici (\ie en utilisant
exclusivement des diff\'erentielles \emph{relatives} et non pas des
diff\'erentielles \emph{absolues}, \ie relatives \`a l'anneau de
constantes absolues $\ZZ$). Comme bien souvent, la consid\'eration
des diff\'erentielles est plus commode ici que celle des
d\'erivations. Notons enfin que si $Y$ est lisse sur~$S$ en $x$, de
dimension relative $n$, alors il existe un voisinage ouvert de~$x$
dans~$Y$ isomorphe \`a un sous-pr\'esch\'ema de
$X=S[t_1,\dots,t_n]$ avec $n=m+1$, comme il r\'esulte de la
d\'efinition et
de~(I~\ref{I.7.6}).

Soient $A$ un anneau noeth\'erien, $x_i$ ($1 \leq i \leq n$) des
\'el\'ements de~$A$, $J$ l'id\'eal engendr\'e par les
$x_i$. On dit que les $x_i$ forment un \emph{système régulier
de générateurs}
\index{regulier (systeme de generateurs)@r\'egulier (syst\`eme de g\'en\'erateurs)|hyperpage}\index{systeme regulier de generateurs)@syst\`eme r\'egulier de g\'en\'erateurs|hyperpage}de~$J$ si l'homomorphisme surjectif canonique
$$
(A/J)[t_1,\dots,t_n] \to \gr^J(A)
$$
d\'efini par les $x_i$ (o\`u le deuxi\`eme membre d\'esigne
l'anneau gradu\'e associ\'e \`a $A$ filtr\'e par les
puissances de~$J$) est un \emph{isomorphisme}. Cette condition
signifie aussi que
\begin{enumerate}
\item[(i)] L'homomorphisme surjectif canonique
$$
S_{A/J} (J/J^2) \to \gr^J(A)
$$
(o\`u, le premier membre d\'esigne l'alg\`ebre sym\'etrique du
$A/J$-module $J/J^2$) est un isomorphisme, et
\item[(ii)] $J/J^2$ est libre et admet pour base les classes des $x_i
\mod J^2$.
\end{enumerate}
Sous cette forme, on voit que si $J\neq A$, les $x_i$ forment un
\emph{système minimal de générateurs} de~$J$, et que
\emph{tout autre système minimal de générateurs} de~$J$
\emph{est un système régulier de générateurs}
(N.B. \og minimal\fg est pris au sens strict: nombre minimum
d'\'el\'ements, qui n'est \'equivalent au sens: minimal pour
l'inclusion, que si $A$ est local); d'autre part, si $J=A$, tout
syst\`eme de g\'en\'erateurs de~$J$ est r\'egulier.

La
condition de r\'egularit\'e d'un syst\`eme de
g\'en\'erateurs d'un id\'eal est stable par localisation par un
ensemble multiplicativement stable quelconque, et d'autre part on voit
tout de suite que pour que $(x_i)$ soit un syst\`eme minimal de
g\'en\'erateurs de~$J$, il suffit d\'ej\`a que pour tout id\'eal
\emph{maximal $\goth m$ contenant $J$}, les $x_i$ d\'efinissent un
syst\`eme r\'egulier de g\'en\'erateurs de~$JA_{\goth m}$ dans
$A_{\goth m}$. Cela nous ram\`ene donc au cas o\`u $A$ est un
anneau local d'id\'eal maximal $\goth m$, et o\`u les $x_i$ sont
dans $\goth m$. \emph{Alors les $x_i$ forment un système
régulier de générateurs de~$J$ si et seulement si ils
forment une $A$-suite au sens de Serre}\footnote{Nous dirons
maintenant plut\^ot \og suite $A$-r\'eguli\`ere\fg, \cf EGA
$0_{\textup{IV}}$~15.1.7 et 15.1.11.}, \ie si pour tout $i$ tel que $1
\leq i \leq n$, $x_i$ est non-diviseur de~$0$ dans
$A/(x_1,\dots,x_{i-1})A$.

Enfin, dans le cas o\`u $A$ est une alg\`ebre sur un anneau $B$,
et o\`u $A/J$ est isomorphe comme $B$-alg\`ebre \`a $B$ (de
sorte que $J$ est le noyau d'un homomorphisme de~$B$-alg\`ebres
$A\to B$), alors les $x_i$ forment un syst\`eme r\'egulier de
g\'en\'erateurs de~$J$ si et seulement si l'homomorphisme
canonique
$$
B[[ t_1,\dots,t_n]]\to \hat{A}
$$
d\'efini par les $x_i$ (o\`u le deuxi\`eme membre d\'esigne le
compl\'et\'e s\'epar\'e $\varprojlim A/J^{n+1}$ de~$A$ pour la
topologie d\'efinie par les puissances de~$J$) est un
\emph{isomorphisme} (il est en tout cas \emph{surjectif}).

Tous ces faits sont bien connus, et figurent sans doute dans le cours
de Serre d'alg\`ebre commutative r\'edig\'e par Gabriel, \`a
peu de choses pr\`es. (O\`u on trouve $N$ autres
caract\'erisations des $A$-suites, dans le cas o\`u $A$ est un
anneau local).

Soit $J$ un id\'eal dans un anneau noeth\'erien $A$. On dira que
$J$ est un \emph{idéal régulier}
\index{ideal regulier@id\'eal r\'egulier|hyperpage}\index{regulier (ideal)@r\'egulier (id\'eal)|hyperpage}si pour tout id\'eal premier $\goth p$ de~$A$, $JA_{\goth p}$ admet
un syst\`eme r\'egulier de g\'en\'erateurs. Il suffit
\'evidemment de le v\'erifier pour $\goth p \supset J$, et on peut
de plus se borner \`a $\goth p$ maximal. Plus g\'en\'eralement,
soit $\cal{J}$ un Id\'eal sur un pr\'esch\'ema localement
noeth\'erien $X$, on dit que $\cal{J}$ est un \emph{Idéal
régulier} si pour tout $x\in X$, $\cal{J}_x$ est un id\'eal de
$\cal{O}_x$ qui admet un syst\`eme r\'egulier de
g\'en\'erateurs. Cela \'equivaut \`a la conjonction des deux
conditions suivantes:
\begin{enumerate}
\item[(a)] L'homomorphisme canonique surjectif
$$
S_{\cal{O}_X/\cal{J}}(\cal{J}/\cal{J}^2) \to \gr^{\cal{J}}(\cal{O}_X)
$$
est un isomorphisme et
\item[(b)] Le faisceau de~$\cal{O}_X/\cal{J}$-Modules
$\cal{J}/\cal{J}^2$ est localement libre.
\end{enumerate}

On
dit alors aussi que le sous-pr\'esch\'ema $Y$ de~$X$ d\'efini
par $\cal{J}$ (donc tel que $\cal{O}_Y$ prolong\'e par $0$ soit
isomorphe \`a $\cal{O}_X/\cal{J}$) est \emph{régulièrement
immergé}
\index{regulierement immerge (preschema)@r\'eguli\`erement immerg\'e (pr\'esch\'ema)|hyperpage}dans~$X$, et on d\'efinit de m\^eme (de fa\c con
\'evidente) la notion de morphisme d'\emph{immersion
régulière},
\index{immersion r\'eguli\`ere|hyperpage}\index{reguliere (immersion)@r\'eguli\`ere (immersion)|hyperpage}(\resp \emph{régulière en un point $x$}), morphisme
d'immersion $Y\to X$ identifiant $Y$ (\resp un voisinage convenable de
$x$), \`a un sous-pr\'esch\'ema ferm\'e r\'eguli\`erement
immerg\'e dans un ouvert de~$X$. (Il ne faut pas dire:
sous-pr\'esch\'ema r\'egulier, car cela signifierait que les
anneaux locaux de~$Y$ sont r\'eguliers). Enfin, des sections $x_i$
de~$\cal{J}$ sont appel\'ees \emph{système régulier de
générateurs} si pour tout $x\in X$, les \'el\'ements
correspondants de~$\cal{O}_x$ forment un syst\`eme r\'egulier de
g\'en\'erateurs de~$\cal{J}_x$ (terminologie compatible avec celle
introduite pour des g\'en\'erateurs d'un id\'eal d'un
anneau). Cela signifie aussi que l'homomorphisme surjectif canonique
$$
\cal{O}_Y[t_1,\dots,t_n] \to \gr^{\cal{J}} (\cal{O}_X)
$$
d\'efini par les $x_i$ est un isomorphisme. Si on sait par avance
que l'Id\'eal $\cal{J}$ est r\'egulier, cela signifie aussi,
simplement, que en tout point $x$ \emph{de~$Y$}, les $x_i$
d\'efinissent une \emph{base} de~$\cal{J}/\cal{J}^2$ sur~$\cal{O}_{Y,x}$. (N.B. cette condition est vide si $Y$ est
vide). Ainsi, pour que $\cal{J}$ admette un syst\`eme r\'egulier
de g\'en\'erateurs, il faut et il suffit que $\cal{J}$ soit
r\'egulier, et le $\cal{O}_Y$-Module $\cal{J} /\cal{J}^2$ soit
globalement libre (et non-seulement localement libre), \ie que
l'homomorphisme canonique $S_{\cal{O}_Y}(\cal{J}/\cal{J}^2) \to
\gr^{\cal{J}} (\cal{O}_X)$ soit surjectif, et que le
$\cal{O}_Y$-Module $\cal{J}/\cal{J}^2$ soit globalement libre.

Un \emph{anneau augmenté est dit régulier}
\index{anneau augment\'e r\'egulier|hyperpage}\index{regulier (anneau augmente@r\'egulier (anneau augment\'e)|hyperpage}si l'id\'eal de l'augmentation est r\'egulier. Ainsi, si $A$ est
un anneau local, consid\'er\'e comme augment\'e dans son corps
r\'esiduel $k$, alors $A$ est un anneau local r\'egulier si et
seulement si c'est un anneau augment\'e r\'egulier.

(\`A vrai dire, il semble qu'il \'etait inutile de commencer par
faire le sorite pr\'eliminaire pour les anneaux, il y a
int\'er\^et \`a commencer avec les faisceaux tout de suite. Si
on veut quelque chose dans le cas noeth\'erien, c'est la
d\'efinition adopt\'ee ici --- a priori moins stricte que celle par
les $A$-suites de Serre --- qui semble pr\'ef\'erable pour les
besoins du calcul diff\'erentiel. Bien entendu, pour bien faire, il
faudrait d\'evelopper aussi au moins une partie de la th\'eorie
des morphismes lisses dans le cadre non-noeth\'erien\footnote{Comme
il est dit dans l'avant-propos, c'est chose faite maintenant, \cf EGA
IV 17, 18}, probablement en
partant du crit\`ere jacobien, de fa\c con \`a obtenir si
possible toutes les propri\'et\'es formelles essentielles des
morphismes lisses et des morphismes \'etales \ie lisses et
quasi-finis; les r\'eciproques seules faisant appel \`a des
hypoth\`eses noeth\'eriennes).
\end{remarques}

Apr\`es ces longs pr\'eliminaires terminologiques, un petit
th\'eor\`eme:
\begin{theoreme}
\label{II.4.15}
Soient $X$ un $S$-pr\'esch\'ema localement de type fini, $Y$ un
sous-pr\'esch\'ema ferm\'e de~$X$ d\'efini par un faisceau
coh\'erent $\cal{J}$ d'id\'eaux sur~$X$, $x$ un point de~$X$. On
suppose maintenant \emph{$Y$ lisse sur~$S$ en $x$} (et rien sur~$X$). Alors les conditions suivantes sont \'equivalentes:
\begin{enumerate}
\item[(i)] $X$ est lisse sur~$S$ en $x$
\item[(ii)] L'immersion $i\colon Y \to X$ est r\'eguli\`ere en
$x$, \ie $\cal{J}_x$ est un id\'eal r\'egulier de~$\cal{O}_x$.
\end{enumerate}
\end{theoreme}

\begin{corollaire}
\label{II.4.16}
Supposons $Y$ \emph{lisse} sur~$S$. Pour que $X$ soit lisse sur~$S$
dans un voisinage de~$Y$ (\ie aux points de~$Y$) il faut et il suffit
que $Y$ soit r\'eguli\`erement plong\'e dans $X$, \ie que
l'immersion $i \colon Y \to X$ soit r\'eguli\`ere.
\end{corollaire}

\subsubsection*{D\'emonstration}
(i) implique (ii). On applique \ref{II.4.10} crit\`ere~(ii), comme
$g\colon X_1 \to X$ est \emph{plat}, pour montrer que l'image inverse
par $g$ du sous-pr\'esch\'ema $Y'$ de~$X'$ est
r\'eguli\`erement plong\'e, on est ramen\'e \`a prouver que
$Y'=S[t_{p+1},\dots, t_n]$ est r\'eguli\`erement plong\'e dans
$S[t_1,\dots,t_n]$, ce qui est trivial (les $t_i$ ($1 \leq i \leq
p$) forment un syst\`eme r\'egulier de g\'en\'erateurs de
l'Id\'eal d\'efinissant $Y'$ dans $X'$).

(ii) implique (i). Soit $g_i$ ($1 \leq i \leq p$) un syst\`eme
r\'egulier de g\'en\'erateurs de~$\cal{J}_x$ et soient $g_i$
($p+1 \leq i \leq n$) des \'el\'ements de~$\cal{O}_{X,x}$ tels que
leurs images $g'_i$ dans $\cal{O}_{Y,x}$ d\'efinissent un morphisme
\emph{étale}
$$
Y_1 \to Y'=S[t_{p+1},\dots, t_n]
$$
d'un voisinages $Y_1$ de~$Y$ dans $Y'$. Les $g_i$ ($1 \leq i \leq n$)
proviennent de sections (de m\^eme nom) de~$\cal{O}_X$ sur un
voisinage $X_1$ de~$x$, et on peut supposer $X_1=X$, $Y_1=Y$. On
obtient ainsi un morphisme
$$
g\colon X \to X'=S[t_1,\dots,t_n]
$$
et tout revient \`a montrer que ce morphisme est \emph{étale}
en~$x$. Prenant $X_1$ assez petit,
on peut supposer que les $g_i$ ($1 \leq i \leq p$) forment un
syst\`eme r\'egulier de g\'en\'erateurs de~$\cal{J}$ sur tout
$X$. En particulier, ils engendrent $\cal{J}$, donc le
sous-pr\'esch\'ema $Y$ de~$X$ s'identifie \`a l'image inverse
par $g$ du sous-pr\'esch\'ema $Y'$ de~$X'$. Soit $x'=g(x)$, alors
la fibre de~$X' \to X$ en $x'$ est
\ignorespaces
identique \`a la fibre de~$Y \to Y'$ en $x$, donc est \'etale sur~$\kres(x')$, donc $g$ est \emph{non ramifié} en $x$, reste \`a
prouver que $g$ est \emph{plat} en $x$. Or le gradu\'e associ\'e
\`a $\cal{O}_{X',x'}$ filtr\'e par les puissances de
$\cal{J}'_{x}$ est \emph{libre} sur~$\cal{O}_{Y',x'}$ en tous
degr\'es, d'autre part le gradu\'e associ\'e \`a
$\cal{O}_{X,x}$ filtr\'e par les puissances de
$\cal{J}_x=\cal{J}'_x\cal{O}_{X,x}$ est isomorphe (sous
l'homomorphisme canonique) au produit tensoriel du pr\'ec\'edent
par $\cal{O}_{Y,x}$ (puisque l'un et l'autre anneau sont des anneaux
de polyn\^omes \`a $n-p$ ind\'etermin\'ees, \`a anneau de
constantes $\cal{O}_{Y',x'}$, \resp $\cal{O}_{Y,x}$), enfin sur~$\cal{O}_{X',x'}/\cal{J}'_{x'}= \cal{O}_{Y',x'}$
$\cal{O}_{X,x}/\cal{J}_x =\cal{O}_{Y,x}$ est plat.

D'apr\`es un crit\`ere g\'en\'eral de platitude (valable pour
un homomorphisme local d'anneaux locaux noeth\'eriens $A' \to A$,
$A'$ \'etant muni d'un id\'eal $J' \neq A'$ tel que le gradu\'e
associ\'e soit libre sur~$A'/J'$ en toute dimension) il s'ensuit que
$X$ est plat sur~$X'$ en $x$, cqfd.

\begin{corollaire}
\label{II.4.17}
Soient $X$ un pr\'esch\'ema localement de type fini sur~$Y$, $i$
une section de~$X$ sur~$Y$, $y$ un point de~$Y$, $x=i(y)$, $\cal{J}$
le faisceau d'id\'eaux sur~$X$ d\'efini par le
sous-pr\'esch\'ema $i(Y)$ (que nous supposons ferm\'e pour
simplifier l'\'enonc\'e, condition v\'erifi\'ee si $X$ est un
sch\'ema).

Les conditions suivantes sont \'equivalentes:
\begin{enumerate}
\item[(i)] $X$ est lisse sur~$Y$ en $x$
\item[(ii)] $i$ est une immersion r\'eguli\`ere en $y$
\item[(iii)] La $\cal{O}_y$-alg\`ebre compl\'et\'ee de
$\cal{O}_x$ pour la topologie d\'efinie par les puissances de
$\cal{J}_x$ est isomorphe \`a une alg\`ebre de s\'eries
formelles $\cal{O}_y[[ t_1,\dots,t_n]]$.
\item[(iii~bis)] Il existe un voisinage ouvert $U$ de~$y$ tel que le
faisceau d'alg\`ebres $\varprojlim i^*(\cal{O}_X/\cal{J}^{n+1})$ sur~$\cal{O}_Y$ soit isomorphe \`a un faisceau de la forme
$\cal{O}_Y[[ t_1,\dots, t_n]]$ au-dessus de~$U$.
\item[(iv)] Il existe un voisinage ouvert $U$ de~$y$, et un voisinage
ouvert $V$ de~$x$, et enfin un $Y$-morphisme $g \colon V \to
U[t_1,\dots, t_n]$, tel que $g$ soit \'etale, que
$i$
induise une section de~$V$ sur~$U$, transform\'ee par $g$ en la
section nulle de~$U[t_1,\dots, t_n]$ sur~$U$.
\end{enumerate}
\end{corollaire}

L'\'equivalence
de (i) et (ii) est un cas particulier de
th\'eor\`eme~\ref{II.4.15}, en faisant $Y=S$, l'\'equivalence de
(ii) et (iii) (et moralement de (ii) et (iii~bis)) a \'et\'e
signal\'e avec les \og rappels\fg. quant \`a l'\'equivalence de (i)
et (iv), elle se d\'eduit facilement de
th\'eor\`eme~\ref{II.4.10} (\'equivalence des conditions (i) et
(ii) dudit).

\begin{corollaire}
\label{II.4.18}
Soit $X$ un pr\'esch\'ema lisse au-dessus de~$S$. Alors le
morphisme diagonal
$$
\Delta_{X/S}\colon X\to X\times_S X
$$
est une
\emph{immersion régulière},
ou comme on dit encore, $X$ est \og \emph{différentiablement
lisse}\fg sur~$S$.
\index{diff\'erentiablement lisse|hyperpage}\end{corollaire}


En effet, c'est un cas particulier de corollaire~\ref{II.4.16}, puisque
$X$ et $X\times_S X$ sont tous deux lisses sur~$S$.

\setcounter{subsection}{17}

\begin{remarques}
\label{rem:II.4.18}
Rappelons (I~\ref{I.1}) que si $X$ est un pr\'esch\'ema au-dessus
de~$S$, on introduit les faisceaux quasi-coh\'erents
d'alg\`ebres
$\cal{P}_{X/S}^n=\cal{O}_{X\times_S X}/\cal{I}_X^{n+1}$ sur~$X$,
(o\`u $\cal{I}_X$ d\'esigne le faisceau
d'id\'eaux
qui d\'efinit la diagonale dans $X\times_S X$), consid\'er\'e comme
faisceau de~$\cal{O}_X$-alg\`ebres gr\^ace \`a la premi\`ere
projection $\pr_1\colon X\times_S X\to X$. Les
$\cal{P}_{X/S}^n$ forment un syst\`eme projectif d'Alg\`ebres sur~$X$, dont la limite projective est not\'ee $\cal{P}_{X/S}^{\infty}$
et n'est autre que le faisceau structural du compl\'et\'e formel
de~$X\times_S X$ le long de la diagonale (en supposant maintenant $X$
localement de type fini sur~$S$, donc les $\cal{P}_{X/S}^n$
coh\'erents). Dire que $X$ est diff\'erentiablement lisse sur~$S$
(\ie que le morphisme diagonal $\Delta_{X/S}$ est une immersion
r\'eguli\`ere) signifie aussi que $\cal{P}_{X/S}^{\infty}$ est
r\'egulier, en tant que faisceau d'alg\`ebres augment\'e vers
$\cal{O}_X$, \ie que $\it{\Omega}_{X/S}^1$ est localement libre et
l'homomorphisme surjectif canonique
$$
S_{\cal{O}_X}(\it{\Omega}_{X/S}^1)\to
\mathrm{gr}_*(\cal{P}_{X/S}^{\infty})
$$
est un isomorphisme, ou enfin que tout point de~$X$
a
un voisinage ouvert sur lequel le faisceau
d'alg\`ebres
augment\'ees $\cal{P}_{X/S}^{\infty}$ soit isomorphe \`a un
faisceau $\cal{O}_X[[t_1,\dots,t_n]]$.

Soit $s$ une section de~$X$ sur~$S$, $\cal{J}$ le faisceau
d'id\'eaux sur~$X$ qu'elle d\'efinit (supposant
pour simplifier que $s(S)$ est ferm\'e), on a alors des
isomorphismes canoniques de~$\cal{O}_X$-alg\`ebres augment\'ees:
\begin{equation}
\label{eq:II.4.4}
s^*(\cal{P}_{X/S}^n)=\cal{O}_X/\cal{J}^{n+1}\quoi,\quad
s^*(\cal{P}_{X/S}^{\infty})= \varprojlim_{n}
\cal{O}_X/\cal{J}^{n+1}
\end{equation}
Ces isomorphismes sont fonctoriels dans un sens \'evident par
changement de base, et (compte tenu de ce fait) redonnent une
caract\'erisation des faisceaux d'alg\`ebres $\cal{P}_{X/S}^n$ sur~$S$. Si par exemple $S=\Spec(k)$, $k$ un corps, alors la donn\'ee
d'une section~$s$ de~$X$ sur~$S$ \'equivaut \`a la donn\'ee d'un
point $x$ de~$X$ rationnel sur~$k$, et les formules
pr\'ec\'edentes signifient que l'on a un isomorphisme de
$k$-alg\`ebres
\begin{equation}
\label{eq:II.4.5}
\cal{P}_{X/S}^n(x)=\cal{O}_x/\goth{m}_x^{n+1}
\end{equation}
ce qui justifie le nom: \og \emph{faisceau des parties principales
d'ordre~$n$ sur~$X$ rel. à~$S$}\fg
\index{parties principales d'ordre~$n$ (faisceau des, syst\`eme des)|hyperpage}\label{indnot:ad1}donn\'e \`a~$\cal{P}_{X/S}^n$. On voit de plus sur
\eqref{eq:II.4.4} que \emph{si $X$ est différentiablement lisse
sur~$S$ en tout point de~$s(S)$, alors $X$ est lisse sur~$S$ en tout
point de~$s(S)$}, (corollaire~\ref{II.4.17}) \emph{la réciproque
étant également vraie} (corollaire~\ref{II.4.18}). Compte tenu
de~\ref{II.4.13}, on en conclut facilement que si $X$ est un
$S$-pr\'esch\'ema localement de type fini, \emph{$X$ est lisse sur~$S$ si et seulement si il est plat sur~$S$ et différentiellement
lisse sur~$S$.} (N.B. l'hypoth\`ese de platitude est essentielle,
comme on voit en prenant pour $X$ un sous-pr\'esch\'ema ferm\'e
de~$S$).

Notons encore, \`a titre de rappel, qu'on obtient une
\emph{deuxième structure d'algèbre} sur~$\cal{P}_{X/S}^n$
gr\^ace \`a la projection $\pr_2\colon X\times_S X\to X$,
se d\'eduisant d'ailleurs de la pr\'ec\'edente \`a l'aide de
l'\emph{involution canonique} du faisceau d'anneaux $\cal{P}_{X/S}^n$,
induit par l'automorphisme de sym\'etrie de~$X\times_S X$. On note
par $d_{X/S}^n$
\label{indnot:bb}ou simplement $d^n$, l'homomorphisme de faisceaux d'anneaux
\begin{equation}
\label{eq:II.4.6}
d_{X/S}^n\colon \cal{O}_X\to\cal{P}_{X/S}^n
\end{equation}
qui correspond \`a cette deuxi\`eme structure
d'Alg\`ebre. Compte tenu de l'isomorphisme~\eqref{eq:II.4.4}, cet
homomorphisme transforme une section $f$ de~$\cal{O}_X$ en une section
$d^n(f)$ de~$\cal{P}_{X/S}^n$ dont l'image inverse par une section $s$
de~$X$ sur~$S$ s'identifie \`a l'image canonique de~$f$ dans
$\Gamma(X,\cal{O}_X/\cal{J}^{n+1})$. Cela justifie le nom de
\og \emph{système des parties principales d'ordre $n$ de~$f$}\fg
\index{systeme des parties principales d'ordre n d'une section@syst\`eme des parties principales d'ordre~$n$ d'une section|hyperpage}donn\'e \`a $d^nf$, notamment dans le cas o\`u $S=\Spec(k)$,
envisag\'e dans la formule~\eqref{eq:II.4.5}.

Pour finir, notons que l'homomorphisme \eqref{eq:II.4.6} peut \^etre
consid\'er\'e comme l'\emph{opé\-ra\-teur différentiel
d'ordre}~$\leq n$\footnote{Pour tout ce qui concerne le pr\'esent
alin\'ea, on pourra consulter EGA~IV~16.8 \`a 16.12.}
(relativement au pr\'esch\'ema des constantes $S$)
\emph{universel}
\index{op\'erateur diff\'erentiel d'ordre~$\leq n$ universel|hyperpage}sur~$\cal{O}_X$, en convenant d'appeler op\'erateur diff\'erentiel
d'ordre~$\leq n$
\index{op\'erateur diff\'erentiel d'ordre~$\leq n$|hyperpage}de~$\cal{O}_X$ dans un Module $F$, un homomorphisme de faisceaux $D$
qui se factorise en
$$
D\colon \cal{O}_X\lto{d^n} \cal{P}_{X/S}^n\lto{u} F
$$
o\`u $u$ est un homomorphisme \emph{de~$\cal{O}_X$-Modules},
d'ailleurs uniquement d\'etermin\'e par~$D$. Cette d\'efinition
concorde avec la d\'efinition r\'ecurrente intuitive ($D$ est un
op\'erateur diff\'erentiel d'ordre $\leq n$ si pour toute section
$g$ de~$\cal{O}_X$ sur un ouvert $U$ de~$X$, $f\mto D(fg)-D(f)$ est
un op\'erateur diff\'erentiel d'ordre $\leq n-1$ sur~$U$). Il
s'ensuit que \emph{si $X$ est différentiablement lisse sur~$S$, le
faisceau d'anneaux des opérateurs différentiels de tous ordres
a la structure simple bien connue} en calcul diff\'erentiel sur les
vari\'et\'es diff\'erentiables, et en particulier admet
localement une base de~$\cal{O}_X$-Module
form\'e 
des \emph{puissances divisées} en des op\'erateurs permutables
$\delta/\delta x_i$ $(1\leq i\leq n)$. Si~$S$ est un faisceau de
$\QQ$-alg\`ebres ($\QQ=$ corps des rationnels) il suffit de prendre
les polyn\^omes ordinaires en les $\delta/\delta x_i$. Dans ce cas
d'ailleurs, et tr\`es exceptionnellement, pour que $X$ soit
diff\'erentiablement lisse sur~$S$, il suffit d\'ej\`a que
$\it{\Omega}_{X/S}^1$ soit localement libre.
\end{remarques}

\begin{remarque}
\label{II.4.19}
\index{Cohen-Macaulay (sch\'ema de)|hyperpage}La terminologie \og immersion r\'eguli\`ere\fg, \og id\'eal
r\'egulier\fg, etc. introduite dans ce num\'ero a rencontr\'e
une opposition assez vive et g\'en\'erale (Chevalley, Serre). On a
propos\'e \og id\'eal de Cohen-Macaulay\fg ou \og id\'eal de
Macaulay\fg ou \og id\'eal macaulayen\fg, ce qui moralement obligerait
\`a adopter aussi \og immersion de Cohen-Macaulay\fg ou \og immersion de
Macaulay\fg. Cette terminologie cependant conflicte avec une autre
d\'ej\`a employ\'ee dans de futures r\'edactions du
multiplodoque, o\`u un morphisme (de type fini) est dit
\og Cohen-Macaulay\fg en un point s'il est plat en ce point, et si la
fibre passant par ce point y a un anneau local qui soit un anneau de
Cohen-Macaulay. En attendant de trouver une solution satisfaisante,
nous garderons sous toutes r\'eserves la terminologie introduite
dans ce num\'ero\footnote{C'est celle adopt\'ee dans
EGA~$\mathrm{0}_\mathrm{IV}$~15.1.7.}.
\end{remarque}

\section{Cas d'un corps de base}
\label{II.5}

\begin{proposition}
\label{II.5.1}
Soient $k$ un corps, $X$ un pr\'esch\'ema de type fini sur~$k$,
$x$ un point de~$X$ et $n$ la dimension de~$X$ en~$x$,
$f\colon X\to\Spec k[t_1,\dots,t_n]=Y$
un morphisme, d\'efini par des \'el\'ements
$f_i\in\Gamma(X,\cal{O}_X)$.
Les conditions suivantes sont \'equivalentes (et entra\^inent que
$X$ est lisse sur~$k$ en $x$, et a fortiori r\'egulier en $x$
d'apr\`es~\ref{II.3.1}):
\begin{enumerate}
\item[(i)] $f$ est \'etale en $x$.
\item[(ii)] Les $df_i$ forment une base de~$\it{\Omega}_{X/k}^1$ en $x$.
\item[(iii)] Les $df_i$ engendrent $\it{\Omega}_{X/k}^1$ en $x$.
\end{enumerate}
\end{proposition}

Comme (i) implique que $X$ est lisse sur~$k$ en $x$, l'implication
(i)$\To$(ii) est un cas particulier de~\ref{II.4.8},
(ii)$\To$(iii) est trivial, reste \`a prouver
(iii)$\To$(i). Or si (iii) est v\'erifi\'e, $f$ est net
en $x$ en vertu de lemme~\ref{II.4.1}, donc (rempla\c cant $x$ par
un voisinage ouvert) quasi-fini, donc dominant par raison de
dimensions. Comme $Y$ est r\'egulier, il s'ensuit que $f$ est
\'etale par (I~\ref{I.9.5}~(ii)) ou (I~\ref{I.9.11}).

\begin{corollaire}
\label{II.5.2}
Sous les conditions pr\'eliminaires de~\ref{II.5.1}, supposons que
$\kres(x)$
soit une extension \emph{finie séparable} de~$k$, et
que les $f_i$ ($1\leq i\leq n$) d\'efinissent des \'el\'ements
de~$\goth{m}_x$. Alors les conditions pr\'ec\'edentes
\'equivalent \`a
\begin{enumerate}
\item[(iv)] Les $f_i$ forment un syst\`eme de g\'en\'erateurs de
$\goth{m}_x$ (ou encore: les $f_i$ mod $\goth{m}_x^2$ forment une base
de~$\goth{m}_x/\goth{m}_x^2$
sur~$\kres(x)$).
\end{enumerate}
\end{corollaire}

En effet, (iv)$\To$(iii) en vertu de la suite exacte
\begin{equation}
\label{eq:II.5.1}
\goth{m}_x/\goth{m}_x^2\to\Omega_{\cal{O}_x/k}^1\to
\Omega_{\kres(x)/k}^1\to 0
\end{equation}
et compte tenu de~$\Omega_{\kres(x)/k}^1=0$ puisque $\kres(x)$ est
\'etale sur~$k$. D'autre part (ii) implique (iv), car comme $X$ et
$\Spec(\kres(x))$ sont lisses sur~$k$ en $x$, on peut dans la suite
exacte pr\'ec\'edente mettre un $0$ sur la gauche en vertu de
\ref{II.4.10}~(iv).

\begin{corollaire}
\label{II.5.3}
Soit $x$ un point de~$X$ (de type fini sur~$k$). Si $X$ est lisse sur~$k$ en $x$, alors $\cal{O}_x$ est r\'egulier, la r\'eciproque
\'etant vraie si
$\kres(x)$
est une extension finie s\'eparable
de~$k$.
\end{corollaire}

En
effet, la r\'eciproque r\'esulte de~\ref{II.5.2}, en prenant un
syst\`eme r\'egulier $(f_i)$ de g\'en\'erateurs de
$\goth{m}_x$. (N.B. au lieu de~\ref{II.5.2}, on peut aussi invoquer le
th\'eor\`eme~\ref{II.4.15}). On conclut:

\begin{proposition}
\label{II.5.4}
Soit $X$ un pr\'esch\'ema de type fini sur~$k$. Si $X$ est lisse
sur~$k$, il est r\'egulier, la r\'eciproque \'etant vraie si $k$
est parfait.
\end{proposition}

Pour la r\'eciproque, on note qu'en vertu de~\ref{II.5.3}, $X$ est
lisse sur~$k$ en tout point ferm\'e, donc partout puisque l'ensemble
des points o\`u il est lisse est ouvert.

\begin{theoreme}
\label{II.5.5}
Soient $X$ un pr\'esch\'ema de type fini sur~$k$, $x$ un point de
$X$, $n$ la dimension de~$X$ en $x$, $k'$ une extension parfaite de
$k$. Les conditions suivantes sont \'equivalentes:
\begin{enumerate}
\item[(i)] $X$ est lisse sur~$k$ en~$x$.
\item[(ii)] $\it{\Omega}_{X/k}^1$ est libre de rang $n$ en~$x$.
\item[(ii~bis)] $\it{\Omega}_{X/k}^1$ est engendr\'e par $n$
\'el\'ements en~$x$.
\item[(iii)] $X$ est diff\'erentiablement lisse sur~$k$ en $x$.
\item[(iv)] Il existe un voisinage ouvert $U$ de~$x$ tel que
$U\otimes_k k'$ soit r\'egulier (\ie les anneaux locaux de ses
points sont r\'eguliers).
\end{enumerate}
\end{theoreme}

On a (i)$\To$(ii) par \ref{II.4.3}~(ii), (ii)$\To$(ii~bis) trivialement et (ii~bis)$\To$(i) gr\^ace
\`a~\ref{II.5.1}. Comme $X$ est plat sur~$k$, on a
(i)$\Leftrightarrow$(iii) en vertu de~\ref{II.4.18}. On a
(i)$\To$(iv) puisque lisse est invariant par extension de la
base et implique r\'egulier, et (iv)$\To$(i) car par
Proposition~\ref{II.5.4}, on voit que $U\otimes_k k'$ est simple sur~$k'$, donc $U$ est simple sur~$k$ par~\ref{II.4.13}.

Prenant pour $x$ le point g\'en\'erique de~$X$ suppos\'e
irr\'eductible, on trouve:

\begin{corollaire}
\label{II.5.6}
Soit $K$ un anneau d'Artin local localis\'e d'une alg\`ebre de
type fini sur le corps~$k$ (par exemple, $K$ peut \^etre une
extension de type fini de~$k$), soit~$n$ le degr\'e de transcendance
sur~$K$ de son corps r\'esiduel. Conditions \'equivalentes:
\begin{enumerate}
\item[(i)] $K$ est une extension finie s\'eparable d'une extension
transcendante pure $k(t_1,\dots,t_n)$ de~$k$.
\item[(ii)]
$\it{\Omega}^1_{X/k}$
est un $K$-module libre de rang~$n$.
\item[(ii~bis)]
$\it{\Omega}^1_{X/k}$
est un $K$-module admettant $n$
g\'en\'erateurs.
\item[(iii)] Le compl\'et\'e $O'$ de~$K\otimes_k K$ pour la
topologie d\'efinie par les puissances de l'id\'eal d'augmentation
$K\otimes_k K\to K$ est une
$K$-alg\`ebre
augment\'ee \og r\'eguli\`ere\fg, \ie isomorphe \`a une
alg\`ebre de s\'eries formelles en~$K$ (N.B. Si~$K$ est un corps,
cela \'equivaut \`a dire que~$O'$ est un anneau local
r\'egulier).
\item[(iv)] $K$ est une extension s\'eparable de~$k$.
\end{enumerate}
\end{corollaire}

En effet, on peut toujours consid\'erer $K$ comme l'anneau local du
point g\'en\'erique d'un sch\'ema de type fini
irr\'eductible~$X$ sur~$k$, et les conditions envisag\'ees sont
les conditions de m\^eme nom dans~\ref{II.5.5}, en prenant dans~(iv)
pour~$k'$ une extension alg\'ebriquement close de~$k$. Seule
l'implication $K$ s\'eparable sur~$k$ $\To$ $X$ lisse sur~$k$ en~$x$, demande une d\'emonstration. Or on est aussit\^ot
ramen\'e gr\^ace \`a~\ref{II.4.13} au cas o\`u le corps de
base est~$k'$, donc alg\'ebriquement clos, donc o\`u il existe un
point $a$ de~$X$ rationnel sur~$k$. Mais alors $X$ est lisse sur~$k$
en~$a$ d'apr\`es~\ref{II.5.4}, a fortiori il est lisse sur~$k$ en le
point g\'en\'erique~$x$, cqfd\footnote{\Cf Errata \`a la fin du pr\'esent Exp.\, II
(p.\, \pageref{II.fin.errata})}.

On remarquera que dans le cas o\`u~$K$ est une extension de type
fini de~$k$, l'\'equivalence de
(i), (ii), (ii~bis), (iv)
est bien
connue, mais que nous ne nous sommes servis d'aucune de
ces
\'equivalences
d\'ej\`a connues. Bien entendu, la
proposition~\ref{II.5.1} contient comme cas particulier qu'une suite
d'\'el\'ements~$x_i$ ($1\leq i\leq n$) est un \og base de
transcendance s\'eparante\fg de~$K$ sur~$k$ si et seulement si
les~$dx_i$ forment une base du $K$-module $\Omega^1_{K/k}$, fait bien
connu par ailleurs.

\begin{corollaire}
\label{II.5.7} Soit $X$ un pr\'esch\'ema de type fini sur un
corps~$k$. Pour que $X$ soit lisse sur~$k$, il faut et il suffit que
$\it{\Omega}^1_{X/k}$
soit localement libre, et que les anneaux locaux des
points g\'en\'eriques des composantes irr\'eductibles de~$X$
soient des extensions s\'eparables de~$k$ (cette derni\`ere
condition \'etant automatiquement v\'erifi\'ee si $k$ est
parfait et $X$ r\'eduit).
\end{corollaire}

On peut supposer $X$ connexe, soit $n$ le rang de
$\it{\Omega}^1_{X/k}$
suppos\'e localement libre. D'apr\`es l'hypoth\`ese
et~\ref{II.5.6}, c'est aussi le degr\'e de transcendance des
extensions
de~$k$ d\'efinies par les anneaux locaux des points
g\'en\'eriques de~$X$, donc toutes les composantes
irr\'eductibles de~$X$ sont de dimension~$n$. On conclut alors
gr\^ace \`a~\ref{II.5.5}.

On fera attention que si $K$ est une extension finie (non
n\'ecessairement s\'eparable) de~$k$, alors $\Omega^1_{K/k}$ est
un $k$-module libre, donc posant $X=\Spec(K)$, $\Omega^1_{X/k}$ est un
faisceau localement libre, et $X$ est r\'eduit, sans que $X$ soit
n\'ecessairement lisse sur~$k$. \'Etendant alors les scalaires \`a
la cl\^oture alg\'ebrique de~$k$, on trouve un exemple analogue,
o\`u $k$ est alg\'ebriquement clos, mais $X$ en revanche
n'\'etant pas r\'eduit.

\begin{corollaire}\label{II.5.8}
Soient $X$ un pr\'esch\'ema de type fini sur le corps~$k$, $x$ un
point de~$X$, $n$ la dimension de~$X$ en~$x$, $p$ la dimension
de~$\cal{O}_x$,
\ie la codimension dans $X$ de l'adh\'erence $Y$ de
$x$ dans~$X$; donc $n-p$ le degr\'e de transcendance
de~$\kres(x)$
sur~$k$. Soient $f_i$ ($1\leq i\leq n$) des \'el\'ements
de~$\cal{O}_x$, tels que $f_i\in \goth{m}_x$ pour $1\leq i\leq p$. Les
conditions suivantes sont \'equivalentes

\begin{enumerate}
\item[(i)] le germe de morphisme en~$x$
$$
X\to \Spec\bigl(k[t_1,\dots,t_n]\bigr)
$$
d\'efini par les $f_i$ est \'etale en~$x$.
\item[(ii)] Les $f_i$ ($1\leq i\leq p$)
engendrent~$\goth{m}_x$,
\ie forment un syst\`eme r\'egulier de param\`etres
de~$\cal{O}_x$, et les classes dans
$\kres(x)$ des~$f_j$ ($p+1\leq j\leq n$) forment une base de transcendance s\'eparante, (\ie les~$d\overline{f}_j$ ($p+1\leq j\leq n$) forment une base
de~$\Omega^1_{\kres(x)/k}$, ou encore engendrent~$\Omega^1_{\kres(x)/k}$).
\end{enumerate}
\end{corollaire}

Supposons~(i) v\'erifi\'e. Il en r\'esulte que les $df_i(x)$
forment une base de~$\it{\Omega}^1_{X/k}(x)$ \eqref{II.4.8} donc leurs
images $d\overline{f}_i(x)$ dans $\Omega^1_{\kres(x)/k}$ engendrent
cet espace vectoriel sur~$k$. Comme les $\overline{f}_i$ pour $1\leq
i\leq p$ sont nuls, il s'ensuit qu'il suffit de prendre les
$d\overline{f}_i(x)$ avec $p+1\leq i\leq n$. Comme le degr\'e de
transcendance de~$\kres(x)$ sur~$k$ est~$n-p$, il r\'esulte alors
du corollaire~\ref{II.5.6} crit\`ere~(iii) (appliqu\'e
\`a~$K=\kres(x)$) que $Y$ est lisse sur~$k$ en son point
g\'en\'erique~$x$, et que les~$d\overline{f}_i(x)$ ($p+1\leq i\leq
n$) forment une \emph{base} de~$\Omega^1_{\kres(x)/k}$
sur~$\kres(x)$. Par suite, la condition~(ii) de~\ref{II.4.9} est
v\'erifi\'ee, donc aussi la condition~(iii) et en particulier
les~$f_i$ ($1\leq i\leq p$) forment un syst\`eme de
g\'en\'erateurs de~$\goth{m}_x$. Comme
$\cal{O}_x$ est de dimension~$p$, ils forment donc un syst\`eme
r\'egulier de param\`etres en~$x$. Cela prouve~(ii).

Supposons~(ii) v\'erifi\'e. En vertu de la suite
exacte~\eqref{II.5.1}, il s'ensuit que les~$df_i(x)$
engendrent~$\it{\Omega}^1_{X/k}$,
d'o\`u~(i) gr\^ace \`a~prop.\, \ref{II.5.1}.

\begin{corollaire}
\label{II.5.9} Soient~$X$ un pr\'esch\'ema de type fini sur le
corps~$k$, $x$ un point de~$X$, $n$ la dimension de~$X$ en~$x$, $p$ la
dimension
de~$\cal{O}_x$,
\ie la codimension de l'adh\'erence~$Y$
de~$x$ dans~$X$, donc~$n-p$ le degr\'e de transcendance
de~$\kres(x)$
sur~$k$. Conditions \'equivalentes:
\begin{enumerate}
\item[(i)] $\cal{O}_x$ est r\'egulier
et~$\kres(x)$
 est une
extension s\'eparable de~$k$.
\item[(ii)] $X$ est lisse sur~$k$ en~$x$, et l'homomorphisme canonique
$$
\goth{m}_x/\goth{m}_x^2\to
\Omega^1_{\cal{O}_x/k}\otimes_{\cal{O}_x}
\kres(x)=\it{\Omega}^1_{X/k}(x)
$$
est injectif.
\item[(iii)] Il y a des~$f_i\in \cal{O}_x$ ($1\leq i\leq n$)
avec~$f_i\in \goth{m}_x$ pour~$1\leq i\leq p$, tels que le germe de
morphisme en~$x$ de~$X$ dans~$\Spec\bigl(k[t_1,\dots,t_n]\bigr)$
d\'efini par les~$f_i$ soit \'etale en~$x$ (\ie par~\ref{II.5.1}
tels que les~$df_i(x)$ engendrent~$\it{\Omega}^1_{X/k}(x)$).
\item[(iv)] Il y a des~$f_i\in \cal{O}_x$ ($1\leq i\leq n$) tels que
les~$f_i$ ($1\leq i\leq p$) engendrent~$\goth{m}_x$ et que
les~$df_j(x)$ ($p+1\leq j\leq n$) 
engendrent~$\Omega^1_{\kres(x)/k}$ sur~$\kres(x)$.
\end{enumerate}
\end{corollaire}

L'\'equivalence de~(iii) et~(iv) r\'esulte du
corollaire~\ref{II.5.8}, ces conditions \'equivalent aussi
d'apr\`es~\ref{II.4.9} au fait que~$X$ est lisse sur~$k$ en~$x$ et
que la condition~(ii)
de~\ref{II.4.10}
est v\'erifi\'ee. Donc elles \'equivalent au fait que~$X$ est
lisse sur~$k$ en~$x$ et que la condition~(iv) de~\ref{II.4.10} est
v\'erifi\'ee, donc \`a~\ref{II.5.9}~(ii). Ou au fait que~$X$ est
lisse sur~$k$ en~$x$ et que la condition~(i) de~\ref{II.4.10} est
v\'erifi\'ee, qui ici signifie simplement
que~$\kres(x)$
est
s\'eparable sur~$k$. Cela implique~\ref{II.5.9}~(i), il reste \`a
prouver que~\ref{II.5.9}~(i) l'implique, \ie \`a prouver le

\begin{corollaire}
\label{II.5.10}
Soit~$x$ un point
d'un
pr\'esch\'ema de type fini sur le corps~$k$, tel
que~$\kres(x)$
soit s\'eparable sur~$k$. Pour que~$X$ soit lisse sur~$k$ en~$x$,
il faut et il suffit qu'il soit r\'egulier en~$x$ (\ie que l'anneau
local~$\cal{O}_x$ soit r\'egulier).
\end{corollaire}

En effet, s'il en est ainsi, on peut \'evidemment trouver
des~$f_i\in \cal{O}_x$ ($1\leq i\leq n$) satisfaisant la
condition~\ref{II.5.9}~(iv).

\subsection*{Errata}
\label{II.fin.errata}
Dans
le pr\'esent num\'ero, d\'emonstration de~\ref{II.5.6}, on
a utilis\'e le fait qu'un sch\'ema de type fini r\'eduit non
vide sur un corps alg\'ebriquement clos admet au moins un point
r\'egulier (donc lisse), fait qui se d\'emontre d'habitude par
voie diff\'erentielle (via le th\'eor\`eme de Zariski que
l'ensemble des points r\'eguliers de~$X$ est ouvert). Si on veut
\'eviter un cercle vicieux, il faut d\'emontrer que si~$K/k$ est
une extension s\'eparable de type fini, et si les~$f_i\in K$ sont
tels que~$d_{K/k}f_i$ forment une base de~$\Omega^1_{K/k}$, ($1\leq
i\leq n$), alors~$n$ est le degr\'e de transcendance de~$K$
sur~$k$,
\ie les~$f_i$ sont alg\'ebriquement ind\'ependants. La
d\'emonstration de ce fait \`a l'aide du crit\`ere de
Mac Lane
est bien connue, \cf Bourbaki, Alg\`ebre, Chap.\, V par.\, 9 th.\, 2: on
prend un polyn\^ome~$g\in k[t_1,\dots,t_n]$ de degr\'e minimal
tel que~$g(f_1,\dots,f_n)=0$, on a donc
$$
\sum\frac{dg}{d t_i}(f_1,\dots,f_n) df_i=0
$$
d'o\`u (puisque les~$df_i$ forment une base de~$\Omega^1_{K/k}$) le
fait que les~$dg/d t_i$ annulent~$(f_1,\dots,f_n)$, donc sont nuls
d'apr\`es le caract\`ere minimal de~$g$, donc si~$k$ est de
caract\'eristique~$0$ on a~$g=0$, et si~$k$ est de
caract\'eristique~$p\neq 0$ on a~$g=h(t_1^p,\dots,t_n^p)$.
Utilisant le crit\`ere de
Mac Lane,
on voit que le
polyn\^ome~$h\in k[t_1,\dots,t_n]$ annule aussi~$(f_1,\dots,f_n)$,
d'o\`u encore~$g=0$ par le caract\`ere minimal de~$g$.

